% ============================================================================
% Dynamic Leadership Vitality Theory (DLVT)
% Corrected working-paper manuscript (skeleton, statements reviewed & verified)
%
% Build:  pdflatex main && bibtex main && pdflatex main && pdflatex main
% Bibliography: ../references.bib (repository root)
%
% Numbering conventions (kept in sync with the companion code and tests):
%   - Equations are numbered per section; the load-bearing anchors are
%     Eq. (3.2) impact, (3.5) vitality ODE, (3.6) capital ODE, (3.7) complexity
%     coupling, (3.14) carrying capacity, (3.15) comparative statics.
%   - Definitions 5 (depletion ratio) and 7 (strategic threshold) are cited by
%     dlvt/model.py and dlvt/analysis.py.
%   - Theorem 1 (existence & uniqueness), Theorem 2 (global asymptotic
%     stability), Proposition 3 (carrying capacity) match the code docstrings.
%   - Appendices are numbered A1, A4-A10.  A2-A3 of earlier drafts were
%     retracted (see the note preceding the appendices).
% ============================================================================
\documentclass[11pt,a4paper]{article}

\usepackage[utf8]{inputenc}
\usepackage[T1]{fontenc}
\usepackage{lmodern}
\usepackage[margin=1.1in]{geometry}
\usepackage{amsmath,amssymb,amsthm}
\usepackage{booktabs}
\usepackage{graphicx}
\usepackage[round]{natbib}
\usepackage[hidelinks]{hyperref}

\graphicspath{{../figures/}}

% --- Theorem-like environments (counters are document-global) ---------------
\newtheorem{theorem}{Theorem}
\newtheorem{lemma}{Lemma}
\newtheorem{proposition}{Proposition}
\newtheorem{corollary}{Corollary}
\theoremstyle{definition}
\newtheorem{definition}{Definition}
\newtheorem{assumption}{Assumption}
\theoremstyle{remark}
\newtheorem{remark}{Remark}

% Equations numbered within sections: (3.1), (3.2), ...
\numberwithin{equation}{section}

% --- Shorthand ---------------------------------------------------------------
\newcommand{\Vmax}{V_{\max}}
\newcommand{\Vstrat}{V_{\mathrm{strategic}}}
\newcommand{\Cmax}{C^{*}_{\max}}
\newcommand{\Ctrap}{C_{\mathrm{trap}}}
\newcommand{\dd}{\mathrm{d}}

\title{Dynamic Leadership Vitality Theory:\\
A Closed-Form Capability-Trap Model of Leader Vitality\\
and Career Capital}
\author{William Bendinelli\thanks{%
Corresponding author: \texttt{wbendinelli@gmail.com}.
Working paper. The companion Python package \texttt{dlvt}, all figure
scripts, and the pinning test suite are available in the accompanying
repository; every numerical claim in this manuscript is reproduced by an
automated test.}}
\date{Working paper --- corrected draft, \today}

\begin{document}

\maketitle

\begin{abstract}
\noindent
We develop Dynamic Leadership Vitality Theory (DLVT), the first closed-form,
individual-level dynamical model in which accumulated career capital
endogenously generates the organizational complexity that drains deployment
vitality, yielding a globally asymptotically stable low-vitality/high-status
attractor---and, under a power-law load kernel, an intervention asymmetry in
which equilibrium vitality is invariant to the scope-coupling parameter
$\beta$ and movable only through recovery-side parameters. The model consists
of two coupled ordinary differential equations for subjective vitality $V$ and
career capital $C$, with complexity $O = O_0 + \beta C^{\eta}$ closing the
feedback loop. We prove that the system admits at most one interior
equilibrium for all positive parameters, that its trace is negative at any
interior equilibrium (no Hopf bifurcation), and that the equilibrium is
globally asymptotically stable on the open set $\{C>0\}$ intersected with a
trapping rectangle (Bendixson--Dulac plus Poincar\'e--Bendixson). In contrast
to the bistable tipping structure of organizational capability traps
\citep{repenning2002capability,rahmandad2016capability}, the DLVT trap is the
\emph{default} destination: a single attractor, not one of two basins. The
theory is micro-founded in effort--recovery theory \citep{meijman1998effort}
and subjective vitality \citep{ryan1997subjective} rather than ego depletion,
whose empirical foundation has not survived replication
\citep{hagger2010ego,carter2015series,vohs2021multisite}, and it adds an
energetic-deployment constraint absent from human-capital theory
\citep{becker1993human}. The calibration is illustrative; the regime label is
governed by the ratio $\mu/\alpha$ relative to an explicit threshold. We
derive six falsifiable propositions with declared scope conditions and outline
a pre-registered field test that could kill the theory's central prediction.
\end{abstract}

\vspace{0.5em}
\noindent\textbf{Keywords:} leadership vitality; career capital; capability
trap; dynamical systems; burnout; formal theory; global stability.

% ============================================================================
% MAIN CHAPTERS
% ============================================================================
% ============================================================================
% Chapter 1 — Introduction
% ============================================================================
\section{Introduction}
\label{sec:introduction}

Senior leaders routinely report that the years in which their formal
authority, reputation, and network---their \emph{career capital}---peaked were
also the years in which their capacity to do the work that built that capital
collapsed. Surveys of executives find majorities considering exits for
well-being reasons \citep{deloitte2022wellbeing}, median CEO tenures have
shortened \citep{conference2023tenure}, and managerial burnout has become a
recognized organizational risk rather than a private misfortune
\citep{maslach2001job,veldhuis2020burnout,muller2026systematic,tu2025dancing}.
Existing accounts treat this pattern either as a resource problem (demands
exceed resources; \citealp{demerouti2001job,bakker2017jdr}), as a stress
process \citep{hobfoll1989conservation}, or as an organizational-level
capability trap \citep{repenning2001nobody,repenning2002capability}. What is
missing is a formal, individual-level theory of \emph{why success itself}
should systematically produce the depletion---and of which interventions can
and cannot reverse it.

This paper supplies that theory. Dynamic Leadership Vitality Theory (DLVT) is
the first closed-form, individual-level dynamical model in which accumulated
career capital endogenously generates the organizational complexity that
drains deployment vitality, yielding a \emph{globally} asymptotically stable
low-vitality/high-status attractor---and, under a power-law load kernel, an
intervention asymmetry in which equilibrium vitality is invariant to the
scope-coupling parameter $\beta$ and movable only through recovery-side
parameters.

The model is deliberately minimal: two coupled ordinary differential
equations. Vitality $V(t)$ recovers toward a ceiling $\Vmax$ and is drained at
a rate that grows nonlinearly with organizational complexity $O(t)$; career
capital $C(t)$ accumulates through energy-gated impact and depreciates; and
complexity is an increasing function of the leader's own capital,
$O = O_0 + \beta C^{\eta}$. The loop closes: capital builds complexity,
complexity drains vitality, vitality gates the impact that builds capital.
The system settles---from essentially every initial condition---into a single
interior equilibrium at which status is high and vitality is chronically low.

\paragraph{Position relative to capability-trap models.}
DLVT is closest in spirit to the system-dynamics literature on capability
traps and capability erosion
\citep{repenning2002capability,repenning2002simulation,rahmandad2016capability,sterman2000business}.
Those models are organizational-level and, crucially, \emph{bistable}: firms
tip into the trap when short-run pressure crosses a threshold, and a
sufficiently large push can tip them back out. DLVT's contrast is sharp. At
the individual level, with the mechanisms modeled here, the low-vitality
equilibrium is not one basin of two---it is the unique attractor of the
positive orthant. The trap is the \emph{default destination}, not an
unfortunate branch. This difference is not rhetorical: it is a falsifiable
structural prediction (no multistability, no hysteresis, no tipping point in
the baseline kernel family; \S\ref{sec:propositions}), and it reverses the
managerial intuition that a one-time ``reset'' (a sabbatical, a
reorganization) can durably relocate the system.

\paragraph{Micro-foundations.}
The vitality construct and its dynamics are grounded in the effort--recovery
model of work strain \citep{meijman1998effort} and the literature on
subjective vitality as a measurable, fluctuating energetic resource
\citep{ryan1997subjective,nix1999revitalization,quinn2012humanenergy}. We
explicitly do \emph{not} build on ego-depletion theory
\citep{baumeister1998ego}: the depletion effect has failed meta-analytic
scrutiny and multisite preregistered replication
\citep{hagger2010ego,carter2015series,vohs2021multisite}, and a theory whose
core state variable rested on it would inherit that fragility. The
effort--recovery mechanism---load during effort, restoration during respite,
and cumulative strain when recovery is structurally insufficient---requires no
depletable willpower substance and is the empirically surviving foundation.

\paragraph{Relation to human-capital theory.}
DLVT does not compete with \citet{becker1993human}; it \emph{adds} to that
framework an energetic-deployment constraint that human-capital theory does
not model. In Becker's account, accumulated capital is deployed at will; in
DLVT, deployment requires vitality, and the accumulation process itself
erodes the vitality needed for deployment. The result is a wedge between
capital owned and capital usable---formalized here as the energy-gated impact
function---that grows with seniority.

\paragraph{Formalism.}
Methodologically, the paper follows the tradition of small, fully analyzed
dynamical models in adjacent fields: ecology
\citep{may1976simple,murray2002mathematical,scheffer2001catastrophic},
organizational simulation \citep{pluchino2010peter}, and computational
psychopathology \citep{cramer2016depression}. Every analytical claim below is
accompanied by a proof or proof sketch, and every numerical claim is pinned by
an automated test in the companion code, so that manuscript and implementation
cannot silently drift apart.

\paragraph{A note on terminology.}
Earlier drafts of this work labeled the low-vitality/high-status equilibrium
the ``zombie-leader equilibrium.'' We retire ``zombie'' as a defined
theoretical construct: it imported connotations the mathematics does not
support and invited misreading of the equilibrium as a discrete type of
person rather than a region of a continuous state space. Throughout this
corrected manuscript the object is called the \emph{low-vitality equilibrium}
(Definition~\ref{def:lowvitality}); the informal label survives only in
legacy function names in the companion code (e.g., \texttt{is\_zombie}).

\paragraph{Contributions.}
The paper makes four contributions.
(i)~\emph{Theory:} a closed-form, individual-level capability-trap model with
a single globally stable attractor, positioned against the bistable
organizational-trap tradition (\S\ref{sec:foundations},
\S\ref{sec:formulation}).
(ii)~\emph{Corrected mathematics:} existence and uniqueness of the interior
equilibrium in the entire positive orthant (Theorem~\ref{thm:uniqueness}),
trace negativity ruling out Hopf bifurcations (Lemma~\ref{lem:trace}), and
global asymptotic stability on $\{C>0\}$ via a corrected trapping rectangle
(Theorem~\ref{thm:gas}); several claims of earlier drafts are explicitly
retracted and replaced (\S\ref{sec:transient}, Appendices~\ref{app:bifurcation}
and \ref{app:gas}).
(iii)~\emph{Intervention asymmetry:} under the power-law load kernel,
equilibrium vitality is invariant to the scope coupling $\beta$
(Corollary~\ref{cor:scope}); scope reduction alone cannot raise equilibrium
vitality, while recovery-side parameters can---stated honestly as a property
of the kernel family, not a universal law.
(iv)~\emph{Falsifiability:} six propositions with declared scope conditions
and an empirical roadmap culminating in a pre-registered $2{\times}2$
scope-versus-recovery field test (\S\ref{sec:propositions},
\S\ref{sec:discussion}).

\paragraph{Roadmap.}
Section~\ref{sec:foundations} develops the theoretical foundations.
Section~\ref{sec:formulation} states the model, the baseline calibration
(Table~\ref{tab:parameters}), and the formal results.
Section~\ref{sec:results} reports the numerical analysis and the proposition
set. Section~\ref{sec:discussion} discusses limitations---including the
illustrative status of the calibration and the stipulated strategic
threshold---and the validation roadmap. Appendices~A1 and A4--A10 contain the
complete proofs and numerical audits.

% ============================================================================
% Chapter 2 — Theoretical Foundations
% ============================================================================
\section{Theoretical Foundations}
\label{sec:foundations}

This section assembles the four theoretical building blocks that the formal
model in \S\ref{sec:formulation} integrates: (i) an energetic state variable
with defensible micro-foundations; (ii) career capital as an accumulating,
depreciating stock; (iii) complexity as an endogenous load generated by the
leader's own accumulation; and (iv) the capability-trap lens that motivates a
dynamical---rather than cross-sectional---formalization.

\subsection{Energetic micro-foundations: effort--recovery, not ego depletion}
\label{sec:foundations-energy}

The vitality construct $V(t)$ is grounded in two literatures. First, the
effort--recovery model \citep{meijman1998effort} holds that work effort
produces load reactions that are reversible if recovery opportunities are
adequate, and cumulative if they are not; strain is a property of the
\emph{effort--recovery cycle}, not of effort alone. This is inherently a
statement about dynamics---a flow of depletion balanced against a flow of
restoration---and it maps directly onto the recovery and drain terms of the
vitality equation. Second, subjective vitality
\citep{ryan1997subjective,ryan2000self,nix1999revitalization} provides a
validated, fluctuating, self-reported measure of the energy available to the
self, distinct from arousal and from affect, that responds to autonomous
versus controlled regulation. Recovery research operationalizes the
restoration side \citep{sonnentag2007recovery,sonnentag2022recovery}, and the
organizational literature on human energy documents its behavioral
consequences \citep{quinn2012humanenergy,owens2016relational}.

We deliberately do \emph{not} rest the theory on ego depletion
\citep{baumeister1998ego}. Although early meta-analytic work reported a
medium-sized depletion effect \citep{hagger2010ego}, subsequent bias-corrected
meta-analyses found the evidence consistent with a true effect near zero
\citep{carter2015series}, and a large multisite preregistered replication
failed to detect the paradigmatic effect \citep{vohs2021multisite}; see also
\citet{inzlicht2016running} for a mechanistic reinterpretation. The
effort--recovery foundation requires only that sustained load without adequate
recovery degrades available energy over days-to-months timescales---a claim
with robust support in occupational health research
\citep{demerouti2001job,bakker2017jdr,alarcon2011meta}---and is therefore the
appropriate micro-foundation for a model whose time unit is weeks, not
laboratory minutes.

\subsection{Career capital as a stock: extending human-capital theory}
\label{sec:foundations-capital}

Career capital $C(t)$ aggregates the human, social, and reputational assets a
leader accumulates \citep{becker1993human,day2015leadership}. Two properties
matter for the formalization. First, capital is built through \emph{impactful
deployment}: visible, effective work compounds reputation and network
position, which is why the accumulation term in \S\ref{sec:formulation} is
proportional to realized impact rather than to effort. Second, capital
depreciates: skills obsolesce, networks decay, reputations fade
\citep{march1991exploration}, motivating a standard exponential depreciation
term.

The departure from \citet{becker1993human} is the deployment side. Classical
human-capital theory prices the stock; it does not model the \emph{energy
required to deploy it}. DLVT adds an energetic-deployment constraint: impact
is gated by vitality, so a leader can be capital-rich and impact-poor. We
emphasize that this is an extension, not a refutation---the model reduces to a
standard accumulation--depreciation stock when the energy gate is held open.

\subsection{Complexity as an endogenous load}
\label{sec:foundations-complexity}

The third block is the observation that accumulation itself generates load.
Rising capital brings broader spans of control, more interfaces, more
matters requiring attention, and steeper coordination costs
\citep{garicano2000hierarchies,ocasio1997towards,hambrick2005executive,zhu2022executive}.
The demands literature has long distinguished the level of demands from the
resources available to meet them
\citep{karasek1979job,demerouti2001job,podsakoff2023challenge}; DLVT's
specific move is to make the demand level an \emph{increasing function of the
leader's own accumulated capital}, $O = O_0 + \beta C^{\eta}$, so that
demands are endogenous to success. This is the ``too-much-of-a-good-thing''
logic \citep{pierce2013tmgt} given a dynamic transmission mechanism: the
benefit (capital) and the cost (complexity load) are coupled through a single
state variable rather than juxtaposed cross-sectionally. A scope condition is
declared here and revisited in \S\ref{sec:discussion}: the model treats
complexity as a deterministic image of the leader's own capital, abstracting
from market, peer, and organizational sources of load.

\subsection{Capability traps and the case for a closed-form model}
\label{sec:foundations-traps}

The capability-trap literature
\citep{repenning2001nobody,repenning2002capability,repenning2002simulation,rahmandad2016capability}
shows how organizations that divert effort from capability maintenance to
short-run output can slide into self-reinforcing decline. Two features
characterize those models: they operate at the organizational level, and their
signature dynamic is \emph{bistability}---a healthy basin and a trapped basin
separated by a tipping threshold, with hysteresis under parameter sweeps.

DLVT imports the trap \emph{mechanism} (success-driven load crowding out
regeneration) but arrives at a structurally different conclusion at the
individual level: a \emph{single} globally asymptotically stable low-vitality
attractor (Theorem~\ref{thm:gas}). There is no healthy basin to defend and no
tipping point to avoid; under the baseline kernel family the trap is where
the dynamics go from essentially all initial conditions. The theoretical
payoff of this contrast is precision about interventions: in a bistable world
one large push can move the system across the separatrix, whereas in DLVT
only \emph{parameter} changes---specifically recovery-side changes
(Corollary~\ref{cor:scope})---relocate the attractor itself.

Finally, the choice of a two-dimensional closed-form model over a richer
simulation follows the tradition of \citet{may1976simple},
\citet{murray2002mathematical}, and \citet{strogatz2015nonlinear} in ecology
and applied dynamics, of \citet{pluchino2010peter} in organizational
modeling, and of \citet{cramer2016depression} in formal psychopathology: with
two state variables the entire phase plane can be characterized
analytically---nullclines, equilibria, trapping sets, and Dulac
certificates---so that every qualitative claim is a theorem rather than a
simulation summary. The cost is austerity; the benefit is that the theory can
be wrong in public, in closed form.

% ============================================================================
% Chapter 3 — Mathematical Formulation
%
% Equation-numbering contract with the companion code (do not renumber):
%   (3.2) impact, (3.5) vitality ODE, (3.6) capital ODE, (3.7) complexity
%   coupling, (3.14) carrying capacity, (3.15) comparative statics.
% Definition 5 = depletion ratio; Definition 7 = strategic threshold.
% All other displayed mathematics in this section must remain unnumbered.
% ============================================================================
\section{Mathematical Formulation}
\label{sec:formulation}

\subsection{State variables}
\label{sec:state}

\begin{definition}[Subjective vitality]
\label{def:vitality}
$V(t)$ denotes the leader's subjective vitality at time $t$: the energetic
resource available for deliberate leadership work
\citep{ryan1997subjective,meijman1998effort}. Vitality is measured in
arbitrary energy units and confined to $[0,\Vmax]$, where $\Vmax$ is the
individual's recovery ceiling.
\end{definition}

\begin{definition}[Career capital]
\label{def:capital}
$C(t)\ge 0$ denotes accumulated career capital: the stock of human, social,
and reputational assets that confer status and opportunity
\citep{becker1993human}.
\end{definition}

The state space is
\begin{equation}
\label{eq:statespace}
(V(t),\,C(t)) \in [0,\Vmax]\times[0,\infty),
\end{equation}
and Appendix~\ref{app:invariance} verifies that the flow defined below never
leaves it (positive invariance). Time is measured in weeks; one unit is of
the order of one working week, though the calibration is illustrative
(\S\ref{sec:calibration}).

\subsection{Functional forms}
\label{sec:forms}

\begin{definition}[Organizational complexity]
\label{def:complexity}
$O(t) \ge O_0 > 0$ denotes the organizational complexity load facing the
leader: the volume of interfaces, coordination demands, and matters requiring
attention \citep{garicano2000hierarchies,hambrick2005executive}. $O_0$ is the
irreducible baseline load; the coupling of $O$ to the leader's own capital is
specified in \eqref{eq:complexity} below.
\end{definition}

\begin{definition}[Energy-gated impact]
\label{def:impact}
Leadership impact is
\begin{equation}
\label{eq:impact}
I \;=\; \frac{C\,V}{1+\varphi O},
\end{equation}
with suppression coefficient $\varphi>0$.
\end{definition}

Impact requires both assets ($C$) and energy ($V$)---the energetic-deployment
constraint absent from classical human-capital theory---and is attenuated by
complexity friction in the denominator: as $O\to\infty$, $I\to 0$ for any
finite $(V,C)$.

Vitality dynamics balance two flows. Recovery follows a saturating
(logistic-type) law toward the ceiling,
\begin{equation}
\label{eq:recovery}
\text{recovery}(V) \;=\; R\Bigl(1-\frac{V}{\Vmax}\Bigr),
\end{equation}
with maximal restoration rate $R$ realized at full depletion
\citep{sonnentag2007recovery}. Drain grows as a power law in complexity and is
regularized near $V=0$ by a smooth barrier,
\begin{equation}
\label{eq:drain}
\text{drain}(V,O) \;=\; \delta\,O^{\gamma}\,\frac{V}{V+\varepsilon},
\qquad \gamma \ge 1,\ \varepsilon>0,
\end{equation}
where $\delta$ is the energetic cost coefficient and $\gamma$ the complexity
exponent. The factor $V/(V+\varepsilon)$ encodes that a fully depleted leader
cannot be drained further (disengagement at exhaustion) and guarantees
positive invariance (Appendix~\ref{app:invariance}); analytical results are
stated for $\varepsilon>0$ with the $\varepsilon\to0$ limit tabulated in
Appendix~\ref{app:epsilon}.

\subsection{The dynamical system}
\label{sec:system}

The DLVT model is the autonomous planar system
\begin{equation}
\label{eq:dVdt}
\frac{\dd V}{\dd t}
 \;=\; R\Bigl(1-\frac{V}{\Vmax}\Bigr)
 \;-\; \delta\,O^{\gamma}\,\frac{V}{V+\varepsilon},
\end{equation}
\begin{equation}
\label{eq:dCdt}
\frac{\dd C}{\dd t} \;=\; \alpha\,I \;-\; \mu\,C,
\end{equation}
where $\alpha$ is the capital accumulation rate and $\mu$ the depreciation
rate, and the loop is closed by the complexity coupling
\begin{equation}
\label{eq:complexity}
O \;=\; O_0 + \beta\,C^{\eta},
\qquad \beta>0,\ \eta\ge 1 .
\end{equation}
Substituting \eqref{eq:impact} into \eqref{eq:dCdt} gives
$\dot C = C\,\bigl(\alpha V/(1+\varphi O)-\mu\bigr)$, so capital dynamics
factor through $C$ itself; this structure carries the scope-absorption result
of \S\ref{sec:equilibrium-lowvitality}.

\begin{definition}[Depletion ratio]
\label{def:depletion}
The depletion ratio is
\begin{equation}
\label{eq:gamma-ratio}
\Gamma(t) \;=\; \frac{\delta\,O(t)^{\gamma}}{R}.
\end{equation}
$\Gamma>1$ indicates that the drain ceiling exceeds the recovery ceiling: no
vitality level near the top of the range can be sustained, and the
quasi-equilibrium of \S\ref{sec:transient} collapses toward the depletion
band $\{V<\varepsilon\}$.
\end{definition}

\begin{assumption}[Scope condition: endogenous load]
\label{ass:scope}
Complexity is a deterministic, memoryless image of the leader's own capital,
\eqref{eq:complexity}. Exogenous load shocks, peer effects, and
organizational restructuring are outside the model's scope; the exogenous
variant $\dot V$ with prescribed $C(t)$ (used in
Proposition~\ref{prop:band}(ii)) is the only departure we analyze.
\end{assumption}

\subsection{Positive invariance}
\label{sec:invariance}

At $V=0$ the drain term vanishes and $\dot V|_{V=0}=R>0$, so $V=0$ is
\emph{repelling}; at $V=\Vmax$, $\dot V<0$; at $C=0$, $\dot C=0$, so the
$C$-axis is invariant and capital cannot become negative. Hence
$[0,\Vmax]\times[0,\infty)$ is forward invariant. The elementary verification
is Appendix~\ref{app:invariance}. Two consequences deserve emphasis. First,
trajectories cannot reach $V=0$: the ``finite-time depletion to zero''
theorem of earlier drafts was false as stated and is replaced by
Proposition~\ref{prop:band}. Second, $\{C=0\}$ being invariant means the
global-stability statement of Theorem~\ref{thm:gas} is necessarily about the
\emph{open} set $\{C>0\}$.

\subsection{Nullcline geometry and transient depletion}
\label{sec:transient}

Setting $\dot V=0$ in \eqref{eq:dVdt} and clearing denominators shows that,
for each fixed complexity level $O$, the vitality nullcline solves the
quadratic
\begin{equation}
\label{eq:vqe-quadratic}
\frac{R}{\Vmax}V^{2}
\;-\;\Bigl(R-\frac{R\varepsilon}{\Vmax}-\delta O^{\gamma}\Bigr)V
\;-\;R\varepsilon \;=\;0 .
\end{equation}

\begin{proposition}[Vitality nullcline: single branch, strictly decreasing]
\label{prop:nullcline}
For every $O>0$ (equivalently every $C\ge0$), the quadratic
\eqref{eq:vqe-quadratic} has exactly one positive root $V_{\mathrm{qe}}(O)$,
and along the nullcline $V_{\mathrm{qe}}$ is strictly decreasing in $C$.
\end{proposition}

\begin{proof}
The product of the roots of \eqref{eq:vqe-quadratic} is
$-\varepsilon \Vmax<0$, so the roots are real of opposite sign: exactly one
positive root. Writing $F(V,C)$ for the right side of \eqref{eq:dVdt},
$\partial F/\partial V = -R/\Vmax - \delta O^{\gamma}\varepsilon/(V+\varepsilon)^{2}<0$
and
$\partial F/\partial C = -\delta\gamma O^{\gamma-1}O'(C)\,V/(V+\varepsilon)<0$
for $V,C>0$; by the implicit function theorem
$\dd V_{\mathrm{qe}}/\dd C = -F_C/F_V<0$.
\end{proof}

The positive root $V_{\mathrm{qe}}(O)$ is the \emph{quasi-equilibrium} of the
vitality equation at frozen complexity. At the baseline calibration with
$O=O_0=1$ it sits near the ceiling, $V_{\mathrm{qe}}\approx 9.93$
($\approx0.99\,\Vmax$); the closed-form approximation
$V_{\mathrm{qe}}\approx R\varepsilon/(\delta O^{\gamma})$ used in earlier
drafts is valid \emph{only} in the strong-drain limit
$\delta O^{\gamma}\gg R$ and is badly wrong at moderate load (at $O=1$ it
returns $15>\Vmax$).

\begin{proposition}[Finite-time band entry; corrected]
\label{prop:band}
Consider the vitality equation \eqref{eq:dVdt} with an exogenous complexity
path $O(t)$ (Assumption~\ref{ass:scope}).
\begin{itemize}
\item[(i)] $V=0$ is repelling: $\dot V|_{V=0}=R>0$. In particular no
trajectory reaches $V=0$ in finite or infinite time, and the ``finite-time
depletion to zero'' claim of earlier drafts is retracted.
\item[(ii)] If $O(t)\ge\Omega$ for all $t\ge t_0$ with
$\Omega^{\gamma} > 2R/\delta$, then any trajectory with
$V(t_0)=V_0\ge\varepsilon$ enters the depletion band
$\{V<\varepsilon\}$ no later than
\begin{equation}
\label{eq:Tstar}
T^{*} \;=\; t_0 \;+\;
\frac{V_0-\varepsilon}{\;\tfrac{1}{2}\delta\Omega^{\gamma}-R\;}.
\end{equation}
\item[(iii)] If instead $O(t)\equiv O$ is constant, $V(t)$ converges
monotonically and exponentially to the positive quasi-equilibrium
$V_{\mathrm{qe}}(O)$, the positive root of \eqref{eq:vqe-quadratic}.
\end{itemize}
\end{proposition}

\begin{proof}[Proof sketch]
(i) is immediate from \eqref{eq:dVdt}. (ii) For $V\ge\varepsilon$ we have
$V/(V+\varepsilon)\ge\tfrac12$, so
$\dot V \le R-\tfrac12\delta\Omega^{\gamma}<0$, a uniformly negative bound;
integrating gives \eqref{eq:Tstar}. (iii) At fixed $O$ the scalar flow has a
single positive rest point by Proposition~\ref{prop:nullcline}, with
$F_V<0$ there, so convergence is exponential at rate $|F_V(V_{\mathrm{qe}})|$.
\end{proof}

Thus sustained overload drives vitality into an $\varepsilon$-neighborhood of
zero in finite time---the phenomenologically relevant statement---while zero
itself is never attained.

\subsection{Equilibrium structure: existence and uniqueness}
\label{sec:equilibrium}

Interior equilibria satisfy $\dot C=0$ with $C>0$, which by
\eqref{eq:impact}--\eqref{eq:dCdt} forces $\alpha V/(1+\varphi O)=\mu$, i.e.\
the \emph{capital nullcline}
\begin{equation}
\label{eq:cnullcline}
V_c(O) \;=\; \frac{\mu}{\alpha}\,\bigl(1+\varphi O\bigr),
\end{equation}
an increasing affine function of $O$. Substituting \eqref{eq:cnullcline} into
the vitality equation defines the reduced residual
\begin{equation}
\label{eq:phi-residual}
\Phi(O) \;=\;
R\Bigl(1-\frac{V_c(O)}{\Vmax}\Bigr)
\;-\;\delta\,O^{\gamma}\,\frac{V_c(O)}{V_c(O)+\varepsilon},
\end{equation}
whose roots with $O>O_0$ (equivalently $C>0$) are exactly the interior
equilibria.

\begin{theorem}[Existence and uniqueness, global]
\label{thm:uniqueness}
For every parameter vector in $\mathbb{R}^{11}_{>0}$, the system
\eqref{eq:dVdt}--\eqref{eq:complexity} has \emph{at most one} interior
equilibrium $(V^{*},C^{*})$ in the entire positive orthant. An interior
equilibrium exists if and only if $\Phi(O_0)>0$, which holds at the baseline
calibration.
\end{theorem}

\begin{proof}
Differentiate \eqref{eq:phi-residual}:
\[
\Phi'(O)
= -\frac{R}{\Vmax}\,V_c'(O)
\;-\;\delta\gamma O^{\gamma-1}\,\frac{V_c(O)}{V_c(O)+\varepsilon}
\;-\;\delta O^{\gamma}\,\frac{\varepsilon\,V_c'(O)}{\bigl(V_c(O)+\varepsilon\bigr)^{2}} .
\]
With $V_c'(O)=\mu\varphi/\alpha>0$, every one of the three terms is strictly
negative for $O>0$: the first because $R,\Vmax>0$; the second because
$V_c>0$; the third because $\varepsilon>0$. Hence $\Phi$ is strictly
decreasing on $(0,\infty)$. Since $O(C)=O_0+\beta C^{\eta}$ is strictly
increasing in $C$, the map $C\mapsto\Phi(O(C))$ is strictly decreasing, so it
has at most one root: at most one interior equilibrium, with no restriction
to any bounded region. For existence: $\Phi$ is continuous, strictly
decreasing, and $\Phi(O)\to-\infty$ as $O\to\infty$ (the drain term diverges
while the recovery term is bounded above), so a root with $O>O_0$ exists iff
$\Phi(O_0)>0$. At baseline, $V_c(O_0)=2.3$ and
$\Phi(O_0)=3(1-0.23)-0.02\cdot(2.3/2.4)\approx 2.29>0$.
\end{proof}

Uniqueness in the \emph{entire} positive orthant is stronger than the
bounded-region statements of earlier drafts and matters for what follows: it
excludes multistability outright, for all positive parameter values, and it
is one of the three ingredients of the global-stability proof
(Appendix~\ref{app:gas}).

\subsection{Baseline calibration}
\label{sec:calibration}

Table~\ref{tab:parameters} reports the baseline calibration used in all
figures and in the companion code (initial capital $C_0=5.0$ throughout).

\begin{table}[t]
\centering
\caption{Baseline calibration (Table~1). \emph{The calibration is
illustrative}: it is chosen for a well-conditioned phase portrait in
dimensionless units and has no empirical anchor. No quantitative claim in
this paper should be read as calibrated to data; qualitative claims are
stated with their parameter dependence. Derived quantities are reported at
$\varepsilon=0.1$; see Appendix~\ref{app:epsilon} for the
$\varepsilon\to0$ extrapolation.}
\label{tab:parameters}
\vspace{0.4em}
\begin{tabular}{clll}
\toprule
Symbol & Code key & Baseline & Interpretation \\
\midrule
$R$            & \texttt{R}     & $3.0$  & vitality recovery rate \\
$\Vmax$        & \texttt{Vmax}  & $10.0$ & maximum vitality capacity \\
$\delta$       & \texttt{delta} & $0.02$ & energetic cost coefficient \\
$\gamma$       & \texttt{gamma} & $2.0$  & complexity exponent in drain \\
$O_0$          & \texttt{O0}    & $1.0$  & baseline organizational complexity \\
$\beta$        & \texttt{beta}  & $0.25$ & capital--complexity coupling \\
$\eta$         & \texttt{eta}   & $1.0$  & capital scaling exponent \\
$\alpha$       & \texttt{alpha} & $0.1$  & capital accumulation rate \\
$\varphi$      & \texttt{phi}   & $0.15$ & complexity--impact suppression \\
$\mu$          & \texttt{mu}    & $0.2$  & capital depreciation rate \\
$\varepsilon$  & \texttt{eps}   & $0.1$  & smooth-barrier regularization \\
\midrule
\multicolumn{4}{l}{\emph{Derived quantities at baseline ($\varepsilon=0.1$)}}\\
$V^{*}$        &                & $4.70$   & equilibrium vitality (\S\ref{sec:equilibrium-lowvitality}) \\
$C^{*}$        &                & $32.03$  & equilibrium capital \\
$O^{*}$        &                & $9.01$   & equilibrium complexity \\
$\lambda_{1,2}$&                & $-0.205\pm0.331i$ & Jacobian eigenvalues (App.~\ref{app:jacobian}) \\
$\Cmax$        &                & $45.0$   & carrying capacity, $\eta=1$ (Prop.~\ref{prop:carrying}) \\
$\Cmax|_{\eta=1.5}$ &           & $12.7$   & carrying capacity at $\eta=1.5$ \\
$\Cmax|_{\eta=2}$   &           & $6.7$    & carrying capacity at $\eta=2$ \\
$\Ctrap$       &                & $102.67$ & trapping bound (Thm.~\ref{thm:gas}); $\Ctrap\neq\Cmax$ \\
\bottomrule
\end{tabular}
\end{table}

Because the calibration is illustrative, we report which single ratio governs
the qualitative regime. Along the capital nullcline \eqref{eq:cnullcline},
equilibrium vitality is proportional to $\mu/\alpha$; the classification of
\S\ref{sec:equilibrium-lowvitality} therefore turns on $\mu/\alpha$ relative
to an explicit threshold. At baseline $\mu/\alpha=2.0$, against a flip value
of $2.163$ at which $V^{*}$ crosses $\Vstrat$---an $8\%$ margin. Across a
$\pm2\times$ log-uniform hypercube around the baseline, the probability that
a stable interior equilibrium classifies as low-vitality is approximately
$0.49$ (companion-code robustness sweep). The baseline is thus a genuinely
knife-edge illustration, and we present it as such.

\subsection{The low-vitality equilibrium and scope absorption}
\label{sec:equilibrium-lowvitality}

\begin{definition}[Strategic vitality threshold]
\label{def:threshold}
The strategic vitality threshold is $\Vstrat=\theta\,\Vmax$ with the
stipulated value $\theta=0.5$: the minimum vitality for sustained strategic
(as opposed to reactive) leadership work. The value of $\theta$ is a
labeling convention, not an estimated quantity; sensitivity of the
classification to $\theta$ is reported in \S\ref{sec:limitations}.
\end{definition}

\begin{definition}[Low-vitality equilibrium; outcome regimes]
\label{def:lowvitality}
A stable interior equilibrium $(V^{*},C^{*})$ is a \emph{low-vitality
equilibrium} if $V^{*}<\Vstrat$. A parameter vector is classified
\emph{sustainable} if its stable interior equilibrium has
$V^{*}\ge\Vstrat$, \emph{low-vitality} if $V^{*}<\Vstrat$, and
\emph{collapse-prone} if no stable interior equilibrium exists (equivalently
$\Cmax=0$ in the sense of Proposition~\ref{prop:carrying}). Earlier drafts
called the low-vitality equilibrium the ``zombie-leader equilibrium''; the
label is retired as a construct and survives only in legacy code names.
\end{definition}

At any interior equilibrium, \eqref{eq:cnullcline} gives
\begin{equation}
\label{eq:vstar}
V^{*} \;=\; \frac{\mu}{\alpha}\bigl(1+\varphi O^{*}\bigr),
\end{equation}
and at the baseline calibration the unique interior equilibrium is
\[
V^{*}=4.7025,\qquad C^{*}=32.0337,\qquad O^{*}=9.0084
\qquad(\varepsilon=0.1),
\]
a low-vitality equilibrium sitting $5.9\%$ below $\Vstrat$; the
$\varepsilon\to0$ closed form gives $V^{*}=4.67994$, $C^{*}=31.7325$
(Appendix~\ref{app:equilibrium}). Convergence is a damped spiral with
eigenvalues $-0.205\pm0.331i$ (Appendix~\ref{app:jacobian}).

\begin{corollary}[Scope absorption; intervention asymmetry]
\label{cor:scope}
In the model family \eqref{eq:dVdt}--\eqref{eq:complexity}, where the drain
kernel is the separable power law $\delta f(O)g(V)$ with $f(O)=O^{\gamma}$
and every appearance of $C$ in the flow is channeled through
$O=O_0+\beta C^{\eta}$, the equilibrium pair $(V^{*},O^{*})$ solves the
$\beta$-free system
\[
R\Bigl(1-\frac{V}{\Vmax}\Bigr)=\delta O^{\gamma}\frac{V}{V+\varepsilon},
\qquad
V=\frac{\mu}{\alpha}(1+\varphi O).
\]
Consequently, on the interior branch: (i) $V^{*}$ and $O^{*}$ are invariant
in $\beta$; (ii) the product $\beta\,C^{*\eta}=O^{*}-O_0$ is conserved. At
baseline ($\eta=1$), $\beta C^{*}=8.008$ at $\varepsilon=0.1$ and exactly
$7.933$ in the $\varepsilon\to0$ closed form---the invariant itself is
$\varepsilon$-dependent and must be quoted with its regularization
(Appendix~\ref{app:epsilon}).
\end{corollary}

\begin{proof}
Both equilibrium conditions are equations in $(V,O)$ alone; $\beta$ enters
the dynamics only through the map $C\mapsto O$. Given the unique root
$(V^{*},O^{*})$ (Theorem~\ref{thm:uniqueness}), inverting
\eqref{eq:complexity} yields $C^{*}=((O^{*}-O_0)/\beta)^{1/\eta}$, whence
$\beta C^{*\eta}=O^{*}-O_0$ independent of $\beta$.
\end{proof}

\begin{remark}[Honest status of Corollary~\ref{cor:scope}]
\label{rem:scope}
Earlier drafts headlined this result as a surprising discovery
(``Lemma~2''); the companion code retains that name. It is more honestly a
\emph{reparameterization invariance of the power-law kernel family}: with
$\eta=1$, $\beta$ multiplies a pure power of $C$ and is absorbed by
rescaling the units of capital, $C\mapsto\beta C$. It is \emph{not} a robust
law of leadership dynamics: under load or drain kernels in which the coupling
parameter also sets a saturation or curvature scale that cannot be absorbed
by rescaling $C$ (exponential and Hill specifications are the canonical
cases), the cancellation fails and $\partial V^{*}/\partial\beta\neq0$
generically (\S\ref{sec:robustness}, Appendix~\ref{app:bifurcation}). The
managerial payload survives within the declared family: \emph{scope reduction
alone cannot raise equilibrium vitality; recovery-side interventions---%
raising $R$, lowering $\mu$, raising $\alpha$---can.}
\end{remark}

\subsection{Stability: local, spectral, and global}
\label{sec:stability}

\begin{lemma}[Trace negativity: no Hopf bifurcation]
\label{lem:trace}
At any interior equilibrium of \eqref{eq:dVdt}--\eqref{eq:complexity}, for
all positive parameters, the Jacobian $J$ satisfies
$\operatorname{tr}J<0$ and $\det J>0$. Hence every interior equilibrium is
locally asymptotically stable, and no Hopf bifurcation occurs anywhere in
parameter space.
\end{lemma}

\begin{proof}[Proof sketch]
$J_{VV}=-R/\Vmax-\delta O^{\gamma}\varepsilon/(V+\varepsilon)^{2}<0$
unconditionally. For $J_{CC}$: the condition $\dot C=0$ at an interior
equilibrium forces $\alpha V/(1+\varphi O)=\mu$, so the first and third
terms of
$J_{CC}=\alpha V/(1+\varphi O)
-\alpha CV\varphi\,O'(C)/(1+\varphi O)^{2}-\mu$
cancel, leaving
$J_{CC}=-\alpha C V\varphi\,O'(C)/(1+\varphi O)^{2}<0$.
Thus $\operatorname{tr}J=J_{VV}+J_{CC}<0$. Moreover
$J_{VC}<0$ and $J_{CV}=\alpha C/(1+\varphi O)>0$, so
$\det J=J_{VV}J_{CC}-J_{VC}J_{CV}>0$. A Hopf bifurcation requires a purely
imaginary conjugate pair, i.e.\ $\operatorname{tr}J=0$, which is impossible;
this is consistent with (and implied more globally by) the Dulac certificate
$B=1/C$ of Theorem~\ref{thm:gas}. Full entries and baseline numerics are in
Appendix~\ref{app:jacobian}.
\end{proof}

\begin{proposition}[Carrying capacity as a flux threshold, not a bifurcation]
\label{prop:carrying}
Define the carrying capacity
\begin{equation}
\label{eq:cmax}
\Cmax \;=\;
\Biggl(\frac{(R/\delta)^{1/\gamma}-O_0}{\beta}\Biggr)^{1/\eta},
\end{equation}
the capital level at which the depletion ratio \eqref{eq:gamma-ratio}
reaches $\Gamma=1$ (evaluated at $V=\Vmax$); $\Cmax=44.99$ at baseline. Then
\begin{equation}
\label{eq:cmax-statics}
\frac{\partial \Cmax}{\partial\beta}<0,\qquad
\frac{\partial \Cmax}{\partial\delta}<0,\qquad
\frac{\partial \Cmax}{\partial R}>0 .
\end{equation}
$\Cmax$ is a \emph{flux threshold}: for $C>\Cmax$ the drain ceiling exceeds
the recovery ceiling and vitality is driven toward the depletion band
regardless of its current level. It is \emph{not} a bifurcation point: by
Lemma~\ref{lem:trace}, $\det J>0$ at every interior equilibrium for all
positive parameters, so no saddle-node (fold) bifurcation exists anywhere,
and no equilibria are created or destroyed at $\Cmax$. Earlier drafts'
description of $\Cmax$ as a ``critical threshold beyond which equilibria
cease to exist'' is retracted (Appendix~\ref{app:carrying}).
\end{proposition}

\begin{theorem}[Global asymptotic stability]
\label{thm:gas}
Let the interior equilibrium exist (Theorem~\ref{thm:uniqueness}) and define
the trapping rectangle $\Omega=[0,\Vmax]\times[0,\Ctrap]$ with
\begin{equation}
\label{eq:ctrap}
\Ctrap^{\eta} \;=\;
\frac{\bigl(\alpha\Vmax/\mu-1\bigr)/\varphi \;-\; O_0}{\beta},
\end{equation}
the unique capital level at which $V_c(O(C))=\Vmax$; $\Ctrap=102.67$ at
baseline. Then $\Omega$ is forward invariant, every trajectory with
$V_0\in[0,\Vmax]$, $C_0>0$ enters $\Omega$, and the interior equilibrium is
globally asymptotically stable on the \emph{open} set
$\{C>0\}\cap\Omega$. The set $\{C=0\}$ is invariant, and the axis
equilibrium $(V\approx9.934,\,C=0)$ is a saddle whose stable manifold lies
in $\{C=0\}$, so it attracts no interior trajectory.
\end{theorem}

\begin{proof}[Proof sketch]
(1)~\emph{Trapping:} on $C=\Ctrap$, $\alpha V/(1+\varphi O)\le
\alpha\Vmax/(1+\varphi O(\Ctrap))=\mu$, so $\dot C\le0$; the remaining edges
are handled by \S\ref{sec:invariance}. (2)~\emph{No closed orbits:} the Dulac
function $B(V,C)=1/C$ gives
$\nabla\!\cdot\!(B\mathbf{F})<0$ on $\{V>0,C>0\}$, term by term, ruling out
closed orbits and orbit graphics in the open quadrant (Bendixson--Dulac).
(3)~\emph{Limit sets:} $\Omega$ is compact; by Poincar\'e--Bendixson the
$\omega$-limit set of any interior trajectory is an equilibrium, and by
Theorem~\ref{thm:uniqueness}, Lemma~\ref{lem:trace}, and the saddle structure
of the axis equilibrium, that equilibrium is the interior one. The complete
argument, including the invariance verification edge by edge, is
Appendix~\ref{app:gas}.
\end{proof}

\begin{remark}[$\Ctrap\neq\Cmax$: the corrected trapping bound]
\label{rem:ctrap}
Earlier drafts used the rectangle $[0,\Vmax]\times[0,\Cmax]$ with
$\Cmax=44.99$. That rectangle \emph{leaks}: at $(V,C)=(\Vmax,\,44.99)$ one
computes $\dot C\approx+6.86>0$, and the trajectory from $(V_0,C_0)=(10,40)$
overshoots to $C\approx46.4$ before returning. The trapping bound
\eqref{eq:ctrap} and the carrying capacity \eqref{eq:cmax} are distinct
objects---a forward-invariance bound versus an energy-flux threshold---and
conflating them invalidated the earlier proof
(Appendix~\ref{app:gas}).
\end{remark}

\begin{remark}[Why Poincar\'e--Bendixson]
No elementary global Lyapunov function appears to exist for this system:
weighted quadratic forms centered at the equilibrium fail because the
cross terms generated by the $O^{\gamma}$ drain and the $1/(1+\varphi O)$
gate cannot be dominated globally. The Dulac-plus-Poincar\'e--Bendixson
route is therefore not a convenience but the appropriate tool for this
planar system.
\end{remark}

\subsection{Robustness of the specification}
\label{sec:robustness}

\paragraph{Linear drain ($\gamma=1$).}
With a linear complexity kernel the low-vitality regime disappears at
baseline: the unique stable equilibrium is \emph{sustainable}, at
$V^{*}\approx8.56$ ($C^{*}\approx83.5$). Nonlinear complexity scaling
($\gamma>1$) is therefore a necessary condition for the low-vitality
\emph{regime}---but not for the existence of an attractor, which persists at
$\gamma=1$. The qualitative theory stands or falls with the convexity of the
load--drain relationship, an empirically checkable premise.

\paragraph{Kernel dependence of scope absorption.}
Corollary~\ref{cor:scope} is a property of the power-law family. Under
exponential or Hill (saturating) kernels the $\beta$-invariance of $V^{*}$
fails, and $|\partial V^{*}/\partial\beta|$ grows with the departure from
separability (Proposition~P4 of \S\ref{sec:propositions};
Appendix~\ref{app:bifurcation}). We flag this boundary explicitly so the
intervention asymmetry is not over-generalized.

% ============================================================================
% Chapter 4 — Analysis and Results
% Figures 1-7 are numbered in this section, in order, to match the companion
% code (dlvt/figures.py) and the test-suite docstrings.
% ============================================================================
\section{Analysis and Results}
\label{sec:results}

This section reports the phase-plane analysis at the baseline calibration,
the comparative statics---including the corrected reading of the
$\beta$-sweep---and the falsifiable proposition set. Every number quoted
here is pinned by an automated test in the companion repository.

\subsection{Baseline dynamics}
\label{sec:results-baseline}

Figure~\ref{fig:temporal} traces the baseline trajectory from
$(V_0,C_0)=(8.0,5.0)$ and shows the model's causal loop unfolding in time.
Three features should be visible. First, an early benign phase: complexity
is low, the depletion ratio $\Gamma\ll1$, and vitality holds near the
ceiling while capital compounds. Second, the turn: as capital accumulates,
complexity---and with it the drain term of \eqref{eq:dVdt}---grows
quadratically ($\gamma=2$), and vitality is pulled down through the
strategic threshold even as capital continues to rise. Third, convergence:
the system spirals into the unique interior equilibrium
$(V^{*},C^{*})=(4.7025,\,32.0337)$ with eigenvalues $-0.205\pm0.331i$
(rotation period $\approx19$ time units, amplitude e-folding time
$\approx4.9$; Appendix~\ref{app:jacobian}). The transient overshoot is
diagnostic rather than incidental: from $(10,40)$ capital peaks near
$C\approx46.4$, above the carrying capacity $\Cmax=44.99$, before
complexity friction pulls it back---visual confirmation that $\Cmax$
bounds sustainable flux, not reachable states (Remark~\ref{rem:ctrap}).

\begin{figure}[t]
\centering
\includegraphics[width=0.95\linewidth]{fig1_temporal_evolution.pdf}
\caption{Temporal evolution of $V(t)$, $C(t)$, $O(t)$, and $\Gamma(t)$ at
the baseline calibration (Table~\ref{tab:parameters}), $(V_0,C_0)=(8,5)$.}
\label{fig:temporal}
\end{figure}

Figure~\ref{fig:scenarios} places the three outcome regimes of
Definition~\ref{def:lowvitality} side by side: three qualitatively
different vitality trajectories from identical initial conditions. In the
sustainable panel (low coupling and low cost, $\delta=0.008$,
$\beta=0.15$) vitality settles above $\Vstrat$; in the baseline panel it
settles below; and in the collapse-prone panel (parameters with
$\Cmax=0$) even $C=0$ implies $\Gamma>1$, and vitality is driven into the
depletion band (Proposition~\ref{prop:band}(ii)). The figure establishes
that the regime label is a property of parameters, not of starting
points---the claim that Propositions P1 and P2 below make testable.

\begin{figure}[t]
\centering
\includegraphics[width=0.95\linewidth]{fig2_three_scenarios.pdf}
\caption{Three outcome regimes: sustainable, low-vitality, collapse-prone.}
\label{fig:scenarios}
\end{figure}

Figure~\ref{fig:phase} compresses the same dynamics into the phase plane.
Two curves organize the portrait: the strictly decreasing vitality
nullcline (Proposition~\ref{prop:nullcline}) and the increasing capital
nullcline \eqref{eq:cnullcline}. Their unique interior crossing
(Theorem~\ref{thm:uniqueness}) is the attractor, and the flow field shows
every plotted trajectory spiraling into it---no second equilibrium, no
closed orbit, no basin boundary anywhere in the frame. This is the
single-attractor geometry that separates DLVT from bistable capability
traps.

\begin{figure}[t]
\centering
\includegraphics[width=0.8\linewidth]{fig3_phase_portrait.pdf}
\caption{Phase portrait with nullclines and the unique interior equilibrium
$(V^{*},C^{*})\approx(4.70,\,32.03)$.}
\label{fig:phase}
\end{figure}

\subsection{Comparative statics and the corrected $\beta$-sweep}
\label{sec:results-statics}

Figure~\ref{fig:bifurcation} sweeps $\beta$ and $R$, and the two panels
should look strikingly different. The $R$-panel behaves as intuition
suggests: equilibrium vitality rises with recovery capacity. The
$\beta$-panel carries the substantive result: $V^{*}(\beta)$ is
\emph{flat}, while $C^{*}(\beta)$ falls along the hyperbola
$C^{*}(\beta)=(O^{*}-O_0)/\beta=8.008/\beta$ (at $\varepsilon=0.1$;
$7.933$ in the $\varepsilon\to0$ limit). By Corollary~\ref{cor:scope},
$(V^{*},O^{*})$ solve a $\beta$-free system, so all $\beta$-dependence is
absorbed into capital. Two consequences follow.

First, there is \emph{no bifurcation in $\beta$} within the power-law
family. An earlier draft reported a critical coupling
$\beta_{\mathrm{crit}}\approx0.1015$ at which the equilibrium appeared to
cross the strategic threshold; that value was a scan-window artifact---the
equilibrium branch $C^{*}(\beta)=8.008/\beta$ exits a legacy search window
$C\le80$ precisely at $\beta\approx0.100$---and it is retracted. The full
forensic account, together with the honest interval-reporting protocol now
implemented in the companion code, is Appendix~\ref{app:bifurcation}.

Second, the flatness of $V^{*}(\beta)$ is the formal content of the
\emph{intervention asymmetry}: reducing the scope coupling reallocates the
equilibrium between capital and complexity but leaves equilibrium vitality
untouched, whereas recovery-side parameters move $V^{*}$ directly.

\begin{figure}[t]
\centering
\includegraphics[width=0.95\linewidth]{fig4_bifurcation.pdf}
\caption{Equilibrium loci under parameter sweeps: $C^{*}$ and $V^{*}$
against $\beta$ and $R$. Note the $\beta$-invariance of $V^{*}$
(Corollary~\ref{cor:scope}) and the conserved product $\beta C^{*}$.}
\label{fig:bifurcation}
\end{figure}

Figure~\ref{fig:impact} makes the deployment wedge visible: it plots
realized impact under DLVT against the counterfactual in which the same
capital is deployed without an energy gate---the implicit assumption of
classical human-capital reasoning \citep{becker1993human}. The two curves
coincide early, while vitality is high and the gate is effectively open,
then separate as complexity-driven drain lowers vitality; near equilibrium
the ungated counterfactual continues to credit the leader with the full
yield of a large capital stock while realized impact plateaus well below
it. The gap between the curves is the model's central object: capital
owned but not usable.

\begin{figure}[t]
\centering
\includegraphics[width=0.9\linewidth]{fig5_impact_comparison.pdf}
\caption{Realized impact under DLVT versus the energy-ungated
human-capital counterfactual.}
\label{fig:impact}
\end{figure}

\subsection{Carrying capacity and the regime map}
\label{sec:results-regimes}

Figure~\ref{fig:heatmap} maps the carrying capacity $\Cmax(\beta,R)$ of
Proposition~\ref{prop:carrying}. The level sets should be read as
comparative statics made visible: $\Cmax$ falls with the coupling $\beta$
and the energetic cost $\delta$, and rises with recovery capacity $R$, per
\eqref{eq:cmax-statics}. We reiterate the corrected interpretation: $\Cmax$
is the flux threshold at which $\Gamma=1$ at full vitality, not a fold
point---$\det J>0$ everywhere (Lemma~\ref{lem:trace}), so no equilibria
are created or destroyed there, and trajectories may transiently exceed
it.

Figure~\ref{fig:regimemap} partitions the $(\beta,\delta)$ plane into the
three regimes of Definition~\ref{def:lowvitality}; the reader should note
how close the baseline point sits to the sustainable/low-vitality
boundary. That proximity is quantified by the regime governor: at fixed
drain-side parameters the label turns on $\mu/\alpha$, and the baseline
$\mu/\alpha=2.0$ sits only $8\%$ below the flip value $2.163$ at which
$V^{*}=\Vstrat$. Across a $\pm2\times$ log-uniform hypercube around
baseline, $P(\text{low-vitality}\mid\text{stable
equilibrium})\approx0.49$: the illustrative calibration lies near the
center of its own regime boundary, which is the honest way to read every
baseline classification in this paper.

\begin{figure}[t]
\centering
\includegraphics[width=0.85\linewidth]{fig6_sensitivity_heatmap.pdf}
\caption{Carrying capacity $\Cmax(\beta,R)$; comparative statics per
Eq.~\eqref{eq:cmax-statics}.}
\label{fig:heatmap}
\end{figure}

\begin{figure}[t]
\centering
\includegraphics[width=0.85\linewidth]{fig7_regime_map.pdf}
\caption{Regime classification in the $(\beta,\delta)$ plane
(Definition~\ref{def:lowvitality}).}
\label{fig:regimemap}
\end{figure}

\subsection{Falsifiable propositions}
\label{sec:propositions}

We close the analysis with the proposition set the theory stakes its claim
on. These are \emph{empirical} propositions, labeled P1--P6 to distinguish
them from the formal propositions of \S\ref{sec:formulation}, and each is
stated with the observation that would falsify it. Before stating them, we
declare where they apply. The boxed scope conditions below are part of the
theory's content, not fine print: they specify the leaders, roles, and
observation regimes for which P1--P6 are claims at all, and they follow
directly from the assumptions and robustness analysis of
\S\ref{sec:formulation}.

\begin{center}
\fbox{\begin{minipage}{0.92\linewidth}
\textbf{Scope conditions for P1--P6.} The propositions apply to individual
leaders in roles where (i) complexity load is predominantly generated by the
leader's own accumulated scope (Assumption~\ref{ass:scope}); (ii) the
load--drain relationship is convex ($\gamma>1$, \S\ref{sec:robustness});
(iii) the drain kernel is well approximated by the separable power-law
family (P3--P4 delimit this boundary); and (iv) observation windows are long
enough for equilibrium behavior ($\gtrsim$ several e-folding times;
$\approx5$ time units at baseline). Outside these conditions the theory is
silent, not merely weakened.
\end{minipage}}
\end{center}

\begin{description}
\item[P1 (Single attractor).] We propose that individual vitality--capital
trajectories converge to a single interior steady state, with
damped-spiral transients. \emph{Falsifier:} persistent oscillation, or
evidence of multistability---distinct long-run outcomes for similar
leaders in similar roles that are attributable to initial conditions
rather than parameters.

\item[P2 (Regime governor).] We propose that the low-vitality regime
occurs if and only if the depreciation--accumulation ratio $\mu/\alpha$ is
small relative to the threshold implied by \eqref{eq:vstar}.
\emph{Falsifier:} regime membership uncorrelated with the
depreciation/accumulation ratio across leaders.

\item[P3 (Scope absorption).] We propose that reducing the scope coupling
$\beta$ alone does not durably raise equilibrium vitality: capital
re-expands until the product $\beta C^{\eta}$ is restored.
\emph{Falsifier:} a durable vitality gain from coupling reduction alone,
with $\beta C^{*}$ \emph{not} conserved.

\item[P4 (Boundary of P3).] We propose that the magnitude of
$\dd V^{*}/\dd\beta$ grows with the non-separability of the empirical
drain kernel, so that where drain saturates in complexity (Hill-type),
scope reduction regains leverage. This proposition marks the edge of P3's
validity rather than extending it. \emph{Falsifier:} scope-reduction
effects unrelated to the curvature of the estimated drain kernel.

\item[P5 (Carrying-capacity statics).] We propose that the maximum
sustainable capital $\Cmax$ falls with the coupling $\beta$ and the
energetic cost $\delta$, and rises with recovery capacity $R$, per
\eqref{eq:cmax-statics}. \emph{Falsifier:} observed sustainable-capital
ceilings that do not move with recovery capacity, or that rise with
coupling.

\item[P6 (Intervention ordering).] We propose that recovery-side
interventions---raising $R$, lowering $\mu$, raising $\alpha$---raise
equilibrium vitality, while scope redesign (lowering $\beta$) does not.
\emph{Falsifier:} a well-identified field comparison in which scope
redesign outperforms recovery support on long-run vitality. This is the
theory's killer test (\S\ref{sec:roadmap}).
\end{description}

% ============================================================================
% Chapter 5 — Discussion
% ============================================================================
\section{Discussion}
\label{sec:discussion}

\subsection{Theoretical contributions}
\label{sec:contributions}

DLVT contributes a closed-form, individual-level counterpart to the
organizational capability-trap literature. Where
\citet{repenning2002capability} and \citet{rahmandad2016capability} derive
bistable tipping structures---two basins, a separatrix, hysteresis---DLVT
shows that when the trap mechanism is routed through a single leader's
energetic budget, the structure collapses to one globally asymptotically
stable attractor: the trap is the default destination
(Theorem~\ref{thm:gas}). The practical difference is stark. In a bistable
world, prevention means staying in the healthy basin and rescue means a
large one-time push. In DLVT's world there is no healthy basin to defend;
only \emph{parameter} changes relocate the attractor, and
Corollary~\ref{cor:scope} says which parameters can do so. Relatedly, the
theory extends---rather than contests---human-capital theory
\citep{becker1993human} by adding the energetic-deployment constraint: the
wedge between capital owned and capital usable is the model's central
object, formalized in the energy-gated impact function \eqref{eq:impact}.

\subsection{Managerial implications}
\label{sec:implications}

Within the declared scope conditions, the model's implications invert a
common instinct. Scope reduction---delegating units, flattening
portfolios, trimming spans---feels like the natural remedy for executive
overload, but within the power-law family it is absorbed: capital
re-expands until the complexity load is restored (P3), and equilibrium
vitality is unchanged. What moves the equilibrium is the recovery side:
protected restoration capacity ($R$), slower capital decay for the same
impact ($\mu$), or higher conversion of impact into capital ($\alpha$),
consistent with evidence that respite interventions produce real but
decaying gains unless structurally maintained \citep{davidson2010sabbatical}
and with the recovery literature's emphasis on ongoing rather than episodic
restoration \citep{sonnentag2022recovery}. The model also cautions against
reading high-status incumbents' visible durability as evidence of
sustainability: at the low-vitality equilibrium, status ($C^{*}$) is at its
lifetime maximum precisely while deployable energy is chronically below
threshold.

\subsection{Limitations}
\label{sec:limitations}

Six limitations bound the theory's claims; we state them as bluntly as the
results.

\paragraph{Deterministic dynamics.} The model has no stochastic forcing.
Real careers experience shocks (health events, market dislocations,
reorganizations) that could produce excursions and regime misclassification
in finite samples; the global-stability result guarantees return to the
attractor but says nothing about excursion costs.

\paragraph{Single leader, closed loop.} Complexity is modeled as a
deterministic image of the leader's own capital
(Assumption~\ref{ass:scope}). This is a declared scope condition, not an
empirical claim: peer competition, organizational design, and market
turbulence all generate load that is orthogonal to own-capital
accumulation. Where exogenous load dominates, the exogenous-complexity
variant (Proposition~\ref{prop:band}) is the right object, and the
endogenous feedback results do not apply.

\paragraph{Illustrative calibration.} Table~\ref{tab:parameters} has no
empirical anchor. The honest summary of the baseline is that it sits near
its own regime boundary: $\mu/\alpha=2.0$ against a flip value of $2.163$
($8\%$ margin), and $P(\text{low-vitality}\mid\text{stable equilibrium})
\approx0.49$ across a $\pm2\times$ log-uniform hypercube. Qualitative
structure (uniqueness, stability, scope absorption) is
calibration-independent; regime labels are not.

\paragraph{Stipulated threshold.} The strategic threshold
$\Vstrat=0.5\,\Vmax$ (Definition~\ref{def:threshold}) is a labeling
convention. The baseline equilibrium sits only $5.9\%$ below it, and any
threshold below $0.47\,\Vmax$ reclassifies the baseline as sustainable.
Classification claims should therefore be read as conditional on $\theta$,
and empirical work should estimate the threshold rather than inherit it.

\paragraph{Identifiability.} The eleven parameters collapse to
approximately six effective dimensionless groups after nondimensionalizing
$V$ by $\Vmax$, $C$ by $\beta^{-1/\eta}$, and time by $\mu^{-1}$; within
the equilibrium predictions there are pronounced ridges---most importantly
in $\mu/\alpha$ (which pins $V^{*}$ via \eqref{eq:vstar}) and $R/\delta$
(which pins the flux threshold)---so cross-sectional equilibrium data alone
cannot identify the full parameter vector. Identification requires
trajectories, or shocks (below).

\paragraph{Level of formal guarantee.} The global-stability proof is
planar machinery (Dulac, Poincar\'e--Bendixson) and does not survive added
state dimensions. Extensions---teams, stochastic load, a second resource
stock---will need different tools, and we make no claims beyond two
dimensions.

\subsection{An empirical validation roadmap}
\label{sec:roadmap}

The theory is built to be tested, and the propositions of
\S\ref{sec:propositions} order the program from cheap to decisive.
First, \emph{measurement and trajectory fitting}: subjective vitality
\citep{ryan1997subjective} and recovery experience
\citep{sonnentag2007recovery} instruments administered longitudinally,
with career capital proxied by scope, compensation, and network centrality,
fitted with continuous-time structural equation models to test the damped
spiral signature and single-attractor convergence (P1). Second,
\emph{exogenous shocks to break the $\mu/\alpha$ degeneracy}: promotions,
reorganizations, and mandated leaves act as natural experiments that
separate depreciation from accumulation, which cross-sectional equilibrium
data cannot (P2, \S\ref{sec:limitations}). Third, the \emph{killer test}: a
pre-registered $2{\times}2$ field experiment crossing scope redesign
(coupling reduction) with recovery support (protected restoration
capacity) over a 12--18~month horizon---long enough for the model's
re-expansion dynamics to complete. DLVT predicts a main effect of recovery
support on equilibrium vitality, no durable main effect of scope redesign,
and conservation of the scope--capital product in the redesign arms (P3,
P6). A durable vitality gain from scope redesign alone, with the product
not conserved, falsifies the theory's central mechanism rather than a
calibration choice.

\subsection{Conclusion}
\label{sec:conclusion}

Leadership research has documented, at length, that success is energetically
expensive \citep{hambrick2005executive,rudolph2020healthy}. DLVT's
contribution is to show that a minimal, fully analyzable model of that
observation---capital builds complexity, complexity drains vitality,
vitality gates impact---delivers strong and falsifiable structure: one
equilibrium, globally stable, low-vitality by default under convex load,
indifferent to scope redesign within its kernel family, and movable only
from the recovery side. The corrected mathematics in this manuscript is
deliberately less dramatic than earlier drafts: no finite-time collapse, no
critical coupling, no fold catastrophe. What remains is, we believe, more
useful: a small set of sharp claims, each with a stated boundary and a
stated way to kill it.


% ============================================================================
% APPENDICES
% ============================================================================
\clearpage
\appendix
\renewcommand{\thesection}{A\arabic{section}}
\setcounter{section}{0}
\setcounter{table}{0}
\renewcommand{\thetable}{A\arabic{table}}
\setcounter{figure}{0}
\renewcommand{\thefigure}{A\arabic{figure}}

\paragraph{Note on appendix numbering.}
Appendix numbering follows the working-draft series used by the companion
code and tests (A1, A4--A10). Appendices A2 and A3 of earlier drafts stated a
``finite-time depletion to zero'' result and an associated asymptotic
approximation; both were retracted because $V=0$ is repelling
($\dot V|_{V=0}=R>0$), and their corrected content is absorbed into
\S\ref{sec:transient} (Proposition~\ref{prop:band}) and
Appendix~\ref{app:equilibrium}. The numbering of the surviving appendices is
retained so that references in the code base remain valid.

% ============================================================================
% Appendix A1 — Positive Invariance of the State Space
% ============================================================================
\section{Positive Invariance of the State Space}
\label{app:invariance}

We verify that the physically meaningful region
$[0,\Vmax]\times[0,\infty)$ is forward invariant under
\eqref{eq:dVdt}--\eqref{eq:complexity}, and record the boundary facts used
throughout the paper. The right-hand side is locally Lipschitz on the
region (the barrier $V/(V+\varepsilon)$ is smooth for $V>-\varepsilon$), so
solutions exist and are unique; invariance follows from the sign of the
vector field on each boundary face.

\paragraph{Face $V=0$ (vitality floor).}
At $V=0$ the barrier factor vanishes, $V/(V+\varepsilon)=0$, so the drain
term of \eqref{eq:dVdt} is identically zero and
\[
\dot V\big|_{V=0} \;=\; R \;>\; 0
\qquad\text{for every } C\ge0 .
\]
The vector field points strictly \emph{into} the region: $V=0$ is repelling.
This is the fact that falsifies the ``finite-time depletion to zero''
theorem of earlier drafts and motivates its replacement,
Proposition~\ref{prop:band}. The companion test
\texttt{test\_dv\_dt\_positive\_at\_V\_zero} pins
$\dot V|_{V=0}=R$ exactly, for several values of $C$.

\paragraph{Face $V=\Vmax$ (vitality ceiling).}
At $V=\Vmax$ the recovery term vanishes and, since $O\ge O_0>0$,
\[
\dot V\big|_{V=\Vmax}
\;=\; -\,\delta\,O^{\gamma}\,\frac{\Vmax}{\Vmax+\varepsilon}
\;<\;0 ,
\]
so trajectories cannot cross the ceiling. (Numerically, integration from
$V_0=\Vmax$ respects this to solver tolerance;
\texttt{test\_simulate\_vitality\_stays\_below\_ceiling}.)

\paragraph{Face $C=0$ (capital axis).}
By \eqref{eq:impact}--\eqref{eq:dCdt}, $\dot C = C\bigl(\alpha
V/(1+\varphi O)-\mu\bigr)$, hence
\[
\dot C\big|_{C=0} \;=\; 0
\qquad\text{for every } V ,
\]
so $\{C=0\}$ is an \emph{invariant set}, not merely a reflecting boundary:
capital can neither become negative nor be created from nothing. Dynamics
restricted to the axis reduce to the scalar vitality equation at
$O=O_0$, whose unique rest point $V\approx9.934$ (the positive root of
\eqref{eq:vqe-quadratic} at $O=O_0$) is analyzed in
Appendix~\ref{app:regularization}, where it is shown to be a saddle of the
planar system.

\paragraph{Consequences.}
(i)~$[0,\Vmax]\times[0,\infty)$ is forward invariant; no clamping is needed
in exact arithmetic (the numerical clamps in the companion code are
solver-safety devices only).
(ii)~Because $\{C=0\}$ is invariant, no interior trajectory can reach it in
finite time, and any global-stability statement must be made on the open
set $\{C>0\}$---as in Theorem~\ref{thm:gas}.
(iii)~The barrier constant $\varepsilon$ trades regularity for a small
displacement of equilibrium quantities; the displacement is quantified in
Appendix~\ref{app:epsilon}.

\setcounter{section}{3}% next section prints as A4
% ============================================================================
% Appendix A4 — Equilibrium Algebra and the eps -> 0 Closed Form
% ============================================================================
\section{Equilibrium Algebra and the $\varepsilon\to0$ Closed Form}
\label{app:equilibrium}

This appendix collects the algebra behind Theorem~\ref{thm:uniqueness} and
Corollary~\ref{cor:scope}: the reduction of the equilibrium problem to a
scalar root-finding problem, the exact quadratic in the
$\varepsilon\to0$ limit, and the baseline numerics at both $\varepsilon=0.1$
and $\varepsilon\to0$.

\subsection{Reduction along the capital nullcline}

Interior equilibria require $\dot C=0$ with $C>0$. Factoring
\eqref{eq:dCdt} as $\dot C=C\bigl(\alpha V/(1+\varphi O)-\mu\bigr)$ yields
the capital nullcline \eqref{eq:cnullcline},
$V_c(O)=(\mu/\alpha)(1+\varphi O)$. Substituting into $\dot V=0$ gives the
scalar residual $\Phi(O)$ of \eqref{eq:phi-residual}, strictly decreasing by
Theorem~\ref{thm:uniqueness}. Two structural bounds follow immediately.

\begin{itemize}
\item \emph{Interior lower bound.} Since $O\ge O_0>0$,
$V^{*}=V_c(O^{*})\ge(\mu/\alpha)(1+\varphi O_0)>\mu/\alpha$. At baseline
$V^{*}\ge\mu/\alpha=2$: no interior equilibrium can sit near the vitality
floor, regardless of $\beta$, $\delta$, or $\gamma$. This bound is the
backbone of the regularization audit (Appendix~\ref{app:regularization}).
\item \emph{Interior upper bound.} $V^{*}<\Vmax$ requires
$O^{*}<\bigl(\alpha\Vmax/\mu-1\bigr)/\varphi$, i.e.\ $C^{*}<\Ctrap$ with
$\Ctrap$ as in \eqref{eq:ctrap}: the interior equilibrium always lies
strictly inside the trapping rectangle of Theorem~\ref{thm:gas}.
\end{itemize}

\subsection{The $\varepsilon\to0$ closed form}

For $V>0$ fixed and $\varepsilon\to0$, the barrier $V/(V+\varepsilon)\to1$
and $\dot V=0$ reduces to $R\bigl(1-V/\Vmax\bigr)=\delta O^{\gamma}$. With
the baseline exponent $\gamma=2$ and $V=V_c(O)$, the equilibrium complexity
solves the quadratic
\begin{equation}
\label{eq:eps0-quadratic}
\delta\,O^{2}
\;+\;\frac{R}{\Vmax}\,\frac{\mu\varphi}{\alpha}\,O
\;+\;R\Bigl(\frac{\mu/\alpha}{\Vmax}-1\Bigr)
\;=\;0 .
\end{equation}
The product of the roots is
$\tfrac{R}{\delta}\bigl(\tfrac{\mu/\alpha}{\Vmax}-1\bigr)$, which is
negative precisely when $\mu/\alpha<\Vmax$; in that regime
\eqref{eq:eps0-quadratic} has exactly one positive root---an independent
confirmation of uniqueness in the limit. At baseline the coefficients are
$0.02\,O^{2}+0.09\,O-2.4=0$, whose positive root and induced equilibrium are
\[
O^{*} = 8.93313,\qquad
V^{*} = \frac{\mu}{\alpha}\bigl(1+\varphi O^{*}\bigr) = 4.67994,\qquad
C^{*} = \frac{O^{*}-O_0}{\beta} = 31.7325 ,
\]
with conserved product $\beta C^{*}=O^{*}-O_0=7.93313$
(Corollary~\ref{cor:scope}).

\subsection{Baseline numerics at $\varepsilon=0.1$}

With the regularization retained, $\Phi(O)=0$ is solved numerically
(Brent's method along the nullcline, as implemented in
\texttt{dlvt.analysis.find\_interior\_equilibria}), giving
\[
O^{*}=9.00843,\qquad V^{*}=4.70253,\qquad C^{*}=32.0337,\qquad
\beta C^{*}=8.00843 .
\]
The $\varepsilon$-shift of every equilibrium quantity is below $1\%$ at
$\varepsilon=0.1$ and is tabulated against the closed form in
Appendix~\ref{app:epsilon}. The companion tests pin both sets of values
(\texttt{test\_regularization\_branch\_interior\_is\_unique\_zombie},
\texttt{test\_bifurcation\_interval\_baseline\_beta\_C\_product\_is\_pinned}).

\subsection{Existence condition and regime threshold}

Existence of the interior equilibrium requires $\Phi(O_0)>0$
(Theorem~\ref{thm:uniqueness}), i.e.\ recovery at the nullcline entry point
must exceed baseline drain; a necessary ingredient is
$V_c(O_0)<\Vmax$, i.e.\ $\mu/\alpha<\Vmax/(1+\varphi O_0)$ ($=8.70$ at
baseline). Within the existence region, \eqref{eq:vstar} makes
$V^{*}$ proportional to $\mu/\alpha$ at fixed $O^{*}$, and solving
$V^{*}(\mu/\alpha)=\Vstrat$ numerically gives the flip value
$\mu/\alpha=2.163$ at $\varepsilon=0.1$ ($2.175$ in the $\varepsilon\to0$
closed form), against the baseline $\mu/\alpha=2.0$---the $8\%$ regime
margin quoted in \S\ref{sec:calibration}.

% ============================================================================
% Appendix A5 — Jacobian Entries, Eigenvalues, and Local Stability
% ============================================================================
\section{Jacobian Entries, Eigenvalues, and Local Stability}
\label{app:jacobian}

Write the system \eqref{eq:dVdt}--\eqref{eq:dCdt} as
$\dot V=F(V,C)$, $\dot C=G(V,C)$ with $O=O(C)$ from \eqref{eq:complexity}
and $O'(C)=\beta\eta C^{\eta-1}$.

\subsection{Entries}

Direct differentiation gives, at any point $(V,C)$ with $C>0$:
\begin{align*}
J_{VV} &= \frac{\partial F}{\partial V}
        = -\frac{R}{\Vmax}
          \;-\;\frac{\delta\,O^{\gamma}\,\varepsilon}{(V+\varepsilon)^{2}}
        \;<\;0 ,\\[2pt]
J_{VC} &= \frac{\partial F}{\partial C}
        = -\,\delta\gamma\,O^{\gamma-1}O'(C)\,\frac{V}{V+\varepsilon}
        \;<\;0 \quad (V>0),\\[2pt]
J_{CV} &= \frac{\partial G}{\partial V}
        = \frac{\alpha C}{1+\varphi O}
        \;>\;0 ,\\[2pt]
J_{CC} &= \frac{\partial G}{\partial C}
        = \frac{\alpha V}{1+\varphi O}
          \;-\;\frac{\alpha C V\varphi\,O'(C)}{(1+\varphi O)^{2}}
          \;-\;\mu .
\end{align*}
These are exactly the expressions implemented in
\texttt{dlvt.analysis.jacobian\_eigenvalues}.

\subsection{Simplification at an interior equilibrium}

At an interior equilibrium, $\dot C=0$ with $C^{*}>0$ forces
$\alpha V^{*}/(1+\varphi O^{*})=\mu$, so the first and last terms of
$J_{CC}$ cancel identically and
\[
J_{CC}\big|_{*}
= -\,\frac{\alpha C^{*}V^{*}\varphi\,O'(C^{*})}{(1+\varphi O^{*})^{2}}
\;<\;0 .
\]
Together with $J_{VV}<0$ this proves $\operatorname{tr}J<0$ at
\emph{every} interior equilibrium for \emph{all} positive parameters
(Lemma~\ref{lem:trace}); since additionally $J_{VC}<0<J_{CV}$,
\[
\det J = J_{VV}J_{CC}-J_{VC}J_{CV}
       = \underbrace{J_{VV}J_{CC}}_{>0}
         \;+\;\underbrace{(-J_{VC})\,J_{CV}}_{>0} \;>\;0 .
\]
Trace negative and determinant positive give local asymptotic stability
unconditionally; trace bounded away from zero rules out Hopf bifurcations
throughout parameter space, and determinant bounded away from zero rules
out saddle-node (fold) bifurcations of interior equilibria---the algebraic
heart of Proposition~\ref{prop:carrying}'s ``flux threshold, not
bifurcation'' correction. Both signs are consistent with the global Dulac
certificate $B=1/C$ (Appendix~\ref{app:gas}), which integrates the same
negativity over the trapping region.

\subsection{Baseline numerics}

At the baseline equilibrium
$(V^{*},C^{*},O^{*})=(4.70253,\,32.0337,\,9.00843)$, $\varepsilon=0.1$:
\[
J=
\begin{pmatrix}
-0.30704 & -0.08821\\
\phantom{-}1.36240 & -0.10218
\end{pmatrix},
\qquad
\operatorname{tr}J=-0.40922,\qquad
\det J=0.15155 .
\]
The discriminant $(\operatorname{tr}J)^{2}-4\det J=-0.4388<0$, so the
eigenvalues are the complex pair
\[
\lambda_{1,2} \;=\; -0.205\pm0.331\,i ,
\]
a \emph{damped spiral}: amplitude e-folding time
$1/0.205\approx4.9$ time units and rotation period
$2\pi/0.331\approx19.0$ time units, matching the overshoot-and-return
signature of Figure~\ref{fig:temporal}. The companion test
\texttt{test\_baseline\_equilibrium\_is\_zombie\_and\_stable} pins the
negative real parts and the nonzero imaginary parts of this pair.

% ============================================================================
% Appendix A6 — Carrying Capacity: Derivation and Flux-Threshold Reading
% ============================================================================
\section{Carrying Capacity: Derivation and Flux-Threshold Clarification}
\label{app:carrying}

\subsection{Derivation of Eq.~\eqref{eq:cmax}}

The depletion ratio \eqref{eq:gamma-ratio} compares the drain ceiling
$\delta O^{\gamma}$ (approached as $V/(V+\varepsilon)\to1$) with the
recovery ceiling $R$. Setting $\Gamma=1$ defines the critical complexity
\[
O_{\max} \;=\; \Bigl(\frac{R}{\delta}\Bigr)^{1/\gamma},
\]
and inverting the load kernel \eqref{eq:complexity} gives the capital level
at which the leader's own accumulation pushes complexity to that point:
\[
\Cmax
= \Biggl(\frac{(R/\delta)^{1/\gamma}-O_0}{\beta}\Biggr)^{1/\eta},
\]
which is Eq.~\eqref{eq:cmax}; if $(R/\delta)^{1/\gamma}\le O_0$ we set
$\Cmax=0$ (baseline load alone is unsustainable: the collapse-prone regime
of Definition~\ref{def:lowvitality}). At baseline,
$O_{\max}=\sqrt{150}=12.247$ and $\Cmax=(12.247-1)/0.25=44.99$. The
comparative statics \eqref{eq:cmax-statics} follow by inspection:
$\Cmax$ is decreasing in $\beta$ and $\delta$, increasing in $R$
(both pinned by \texttt{test\_carrying\_capacity\_monotone\_in\_beta} and
\texttt{\ldots\_in\_R}). The sensitivity to the load exponent is steep
(Table~\ref{tab:cmax-eta}).

\begin{table}[h]
\centering
\caption{Carrying capacity at baseline parameters as the load exponent
$\eta$ varies (Proposition~\ref{prop:carrying}).}
\label{tab:cmax-eta}
\vspace{0.4em}
\begin{tabular}{ccc}
\toprule
$\eta$ & $\Cmax$ & pinned by \\
\midrule
$1.0$ & $44.99$ & \texttt{test\_carrying\_capacity\_baseline\_equals\_45}\\
$1.5$ & $12.65$ & \texttt{test\_carrying\_capacity\_eta1\_5\_paper\_value}\\
$2.0$ & $6.71$  & \texttt{test\_carrying\_capacity\_eta2\_paper\_value}\\
\bottomrule
\end{tabular}
\end{table}

\subsection{What $\Cmax$ is: a flux threshold}

For $C>\Cmax$ we have $\Gamma>1$: the maximal drain exceeds the maximal
recovery, so $\dot V<0$ whenever $V/(V+\varepsilon)>R/(\delta O^{\gamma})$,
and the quasi-equilibrium $V_{\mathrm{qe}}(O)$ of
Proposition~\ref{prop:nullcline} falls into the depletion band
(asymptotically $V_{\mathrm{qe}}\approx R\varepsilon/(\delta O^{\gamma})$
in the strong-drain limit). $\Cmax$ therefore marks the capital level
beyond which no high-vitality operating point exists \emph{even
transiently at full recovery}---an energy-flux statement.

\subsection{What $\Cmax$ is not: a bifurcation}

Earlier drafts described $\Cmax$ as ``the critical threshold beyond which
sustainable equilibria cease to exist.'' That language is retracted, for
two independent reasons.

\begin{enumerate}
\item \emph{No fold anywhere.} By Lemma~\ref{lem:trace} and
Appendix~\ref{app:jacobian}, $\det J>0$ at every interior equilibrium for
all positive parameters. A saddle-node (fold) bifurcation requires
$\det J=0$; hence no interior equilibria are created or destroyed at
$\Cmax$ or anywhere else. The equilibrium count is governed solely by the
sign of $\Phi(O_0)$ (Theorem~\ref{thm:uniqueness}).
\item \emph{Trajectories cross it.} $\Cmax$ is not even a dynamical
barrier: the baseline trajectory from $(V_0,C_0)=(10,40)$ transiently
overshoots to $C\approx46.4>\Cmax=44.99$ before spiraling back
(Remark~\ref{rem:ctrap}). A fortiori it cannot serve as a trapping bound,
which is why the global-stability proof uses the distinct constant
$\Ctrap=102.67$ of Eq.~\eqref{eq:ctrap}
(Appendix~\ref{app:gas}).
\end{enumerate}

The corrected vocabulary is: $\Cmax$ = \emph{carrying capacity} (flux
threshold separating sustainable from collapse-prone \emph{regimes} as
parameters vary); $\Ctrap$ = \emph{trapping bound} (forward-invariance
ceiling used in Theorem~\ref{thm:gas}); and the two differ by more than a
factor of two at baseline.

% ============================================================================
% Appendix A7 — Regularization Audit: No Spurious eps-Branch
% ============================================================================
\section{Regularization Audit: No Spurious $\varepsilon$-Branch}
\label{app:regularization}

The smooth barrier $V/(V+\varepsilon)$ in \eqref{eq:drain} is a
positive-invariance device (Appendix~\ref{app:invariance}). A careful
reviewer may worry that it \emph{creates} equilibria: specifically, a
spurious branch at small positive vitality,
$V\approx R\varepsilon/(\delta O^{\gamma})$, with no counterpart in the
$\varepsilon\to0$ limit. This appendix shows the concern is groundless at
baseline---and, more usefully, makes the ``no spurious branch'' claim
executable, via \texttt{dlvt.analysis.find\_regularization\_branch}, so
that any future parameter change reintroducing such a branch is caught by
the test suite.

\subsection{Axis equilibria: exactly one positive root}

On the invariant axis $\{C=0\}$ (where $\dot C\equiv0$), equilibria of the
vitality equation solve the quadratic \eqref{eq:vqe-quadratic} at $O=O_0$:
\[
\frac{R}{\Vmax}V^{2}
-\Bigl(R-\frac{R\varepsilon}{\Vmax}-\delta O_0^{\gamma}\Bigr)V
-R\varepsilon=0 .
\]
The product of its roots is $-\varepsilon\Vmax<0$ for any $\varepsilon>0$
($=-1$ at baseline), so the quadratic has exactly \emph{one} positive
root---there is no second, near-zero axis equilibrium to begin with. At
baseline the positive root is $V\approx9.934$: the axis equilibrium sits
near the \emph{ceiling}, not near zero (the naive approximation
$R\varepsilon/(\delta O_0^{\gamma})=15$ exceeds $\Vmax$ and is invalid at
low load; see \S\ref{sec:transient}).

\subsection{The axis equilibrium is a saddle}

At $(V,C)=(9.934,\,0)$ the Jacobian is upper triangular, since
$J_{CV}=\alpha C/(1+\varphi O)=0$ at $C=0$. Its eigenvalues are the
diagonal entries
\[
\lambda_1=J_{VV}=-0.300
\quad(\text{vitality relaxation along the axis}),\qquad
\lambda_2=J_{CC}=\frac{\alpha V}{1+\varphi O_0}-\mu=+0.664 ,
\]
of opposite sign: a \emph{saddle}. The stable eigenvector $(1,0)$ spans the
axis, so the saddle's stable manifold is contained in $\{C=0\}$; its
unstable direction points into $\{C>0\}$. Consequently the axis equilibrium
attracts no interior trajectory---the fact required by the
global-stability proof (Appendix~\ref{app:gas}) and pinned by
\texttt{test\_regularization\_branch\_axis\_equilibrium\_is\_saddle}.

\subsection{Interior branch: bounded away from zero}

Interior equilibria lie on the capital nullcline \eqref{eq:cnullcline},
hence satisfy the a priori bound
\[
V^{*} \;=\; \frac{\mu}{\alpha}\bigl(1+\varphi O^{*}\bigr)
\;\ge\;\frac{\mu}{\alpha} \;=\;2.0
\quad\text{at baseline},
\]
independent of $\varepsilon$, $\beta$, $\delta$, and $\gamma$. No interior
equilibrium can approach the vitality floor, so no interior
$\varepsilon$-branch exists either. Enumerating all equilibria at baseline
therefore yields exactly two: the interior low-vitality equilibrium
$(4.7025,\,32.0337)$ (stable spiral) and the axis saddle $(9.934,\,0)$
---and nothing below the near-zero threshold
$\min(\mu/\alpha,1)=1.0$
(\texttt{test\_regularization\_branch\_no\_near\_zero\_at\_baseline}).

\subsection{Summary}

The regularization displaces equilibrium quantities by less than $1\%$ at
$\varepsilon=0.1$ (quantified in Appendix~\ref{app:epsilon}) and creates no
equilibria: the product-of-roots argument caps the axis at one positive
root for \emph{any} $\varepsilon>0$, and the nullcline bound keeps the
interior branch above $\mu/\alpha$. The audit is defensive infrastructure,
not a live concern.

% ============================================================================
% Appendix A8 — Bifurcation-Interval Honesty: the Retracted beta_crit
% ============================================================================
\section{Bifurcation-Interval Honesty: the Retracted
$\beta_{\mathrm{crit}}\approx0.1015$}
\label{app:bifurcation}

Earlier drafts reported a critical coupling
$\beta_{\mathrm{crit}}\approx0.1015$ at which the equilibrium vitality
$V^{*}(\beta)$ appeared to cross the strategic threshold $\Vstrat$,
presented as a bifurcation of the baseline model. The claim was wrong
twice over---once numerically and once structurally---and this appendix
documents both failures and the corrected reporting protocol, because the
failure mode (a scan-window artifact masquerading as a critical point) is
generic enough to be worth recording.

\subsection{The structural impossibility}

By Corollary~\ref{cor:scope}, $V^{*}$ is invariant in $\beta$ on the
interior branch: at baseline it is the constant $4.7025$
($\varepsilon=0.1$) for \emph{every} $\beta$ in the structurally interior
region. A $V^{*}(\beta)$ crossing of $\Vstrat=5$ therefore cannot exist at
this calibration. Independently, Lemma~\ref{lem:trace} gives
$\det J>0$ everywhere, so no fold in $\beta$ (or any other parameter)
exists either. Whatever the legacy sweep detected, it was not a
bifurcation.

\subsection{The numerical forensics}

The legacy equilibrium scan searched for nullcline crossings in the window
$C\in[0.01,\,80]$. But the interior branch obeys
\[
C^{*}(\beta)\;=\;\frac{\beta C^{*}}{\beta}\;=\;\frac{8.008}{\beta}
\qquad(\varepsilon=0.1),
\]
which exceeds the window ceiling $C_{\max}^{\mathrm{scan}}=80$ precisely
when
\[
\beta \;<\; \frac{8.008}{80} \;=\; 0.1001 .
\]
For $\beta$ below this value the scan silently \emph{lost} the interior
equilibrium; the interpolation across the last resolved grid points then
manufactured an apparent threshold crossing at $\beta\approx0.1015$. The
``critical coupling'' was the image of the scan window, not of the model:
false precision produced by an implementation constant.

\subsection{The corrected protocol}

The companion function
\texttt{dlvt.analysis.estimate\_bifurcation\_interval} replaces the point
estimate with two honest questions: (i)~does a $V^{*}$-crossing of
$\Vstrat$ exist anywhere in the sweep, once the scan window is made large
enough ($C\le500$) to hold the low-$\beta$ branch? and (ii)~if a crossing
exists (non-baseline calibrations), what is the \emph{interval} of
numerical estimates across a grid of regularizations
$\varepsilon\in\{0.01,0.05,0.1,0.2\}$ and scan resolutions---the
methods-honest reporting granularity, in place of a four-significant-figure
point. At baseline the answers are: no crossing exists;
$V^{*}(\beta)$ is numerically constant at $4.7025$; and the conserved
product is $\beta C^{*}=8.0083$ (all three pinned by the
\texttt{test\_bifurcation\_interval\_*} group, whose diagnostic string is
required to name the scan-window artifact explicitly so the legacy figure
cannot be reproduced by accident).

\subsection{Where real $\beta$-sensitivity lives}

The retraction is scoped, not global. Outside the separable power-law
family---exponential or Hill (saturating) load/drain kernels in which the
coupling parameter also sets a saturation or curvature scale---the
$\beta$-invariance of Corollary~\ref{cor:scope} fails,
$\partial V^{*}/\partial\beta\neq0$, and genuine threshold crossings in
$\beta$ can occur (Proposition~P4, \S\ref{sec:propositions}). Establishing
the empirical kernel family is therefore a first-order measurement task,
not a modeling nicety: it decides whether scope redesign has any
equilibrium leverage at all.

% ============================================================================
% Appendix A9 — eps-Extrapolation of Equilibrium Quantities
% ============================================================================
\section{$\varepsilon$-Extrapolation of Equilibrium Quantities}
\label{app:epsilon}

Analytical statements in this paper are made for the regularized system
($\varepsilon>0$) with the $\varepsilon\to0$ limit available in closed form
(Appendix~\ref{app:equilibrium}); the companion code integrates at
$\varepsilon=0.1$. Because the scope-absorption invariant
$\beta C^{*\eta}=O^{*}-O_0$ (Corollary~\ref{cor:scope}) is itself
$\varepsilon$-dependent, quantitative uses of the invariant must quote the
regularization at which it was computed. Table~\ref{tab:epsilon} reports
the interior equilibrium across a decreasing $\varepsilon$ ladder, computed
with the production root-finder
(\texttt{find\_interior\_equilibria}, $C$-window $200$, $16{,}000$ scan
points), together with the exact $\varepsilon\to0$ closed form of
Eq.~\eqref{eq:eps0-quadratic}.

\begin{table}[h]
\centering
\caption{Interior equilibrium versus the barrier regularization
$\varepsilon$ at baseline parameters. The bottom row is the closed-form
$\varepsilon\to0$ limit, not an extrapolation.}
\label{tab:epsilon}
\vspace{0.4em}
\begin{tabular}{lcccc}
\toprule
$\varepsilon$ & $V^{*}$ & $C^{*}$ & $O^{*}$ & $\beta C^{*}$ \\
\midrule
$0.2$   & $4.72457$ & $32.32766$ & $9.08191$ & $8.08191$ \\
$0.1$ (baseline) & $4.70253$ & $32.03372$ & $9.00843$ & $8.00843$ \\
$0.05$  & $4.69130$ & $31.88406$ & $8.97102$ & $7.97102$ \\
$0.02$  & $4.68450$ & $31.79338$ & $8.94834$ & $7.94834$ \\
$0.01$  & $4.68222$ & $31.76300$ & $8.94075$ & $7.94075$ \\
\midrule
$0$ (closed form) & $4.67994$ & $31.73254$ & $8.93313$ & $7.93313$ \\
\bottomrule
\end{tabular}
\end{table}

Three observations.

\begin{enumerate}
\item \emph{Monotone, small, first-order.} Every quantity approaches its
limit monotonically from above, with increments approximately proportional
to $\varepsilon$ (halving $\varepsilon$ roughly halves the remaining gap).
The total displacement from the working value $\varepsilon=0.1$ to the
limit is $\Delta V^{*}=0.0226$ ($0.48\%$) and
$\Delta(\beta C^{*})=0.0753$ ($0.94\%$)---the source of the ``$<2\%$
shift'' statement in the companion code documentation, here sharpened.
\item \emph{The invariant's two quotable values.} The scope-absorption
product is $\beta C^{*}=8.008$ at the working regularization
$\varepsilon=0.1$ and exactly $7.933$ ($=O^{*}-O_0$ from the quadratic
\eqref{eq:eps0-quadratic}) in the $\varepsilon\to0$ closed form. Both are
correct; neither should be quoted without its $\varepsilon$. The pinning
tests use the $\varepsilon=0.1$ value.
\item \emph{Classifications are $\varepsilon$-robust here.} The regime
label at baseline (low-vitality: $V^{*}<\Vstrat=5$) is unchanged across
the entire ladder, and the regime flip value of $\mu/\alpha$ moves only
from $2.163$ ($\varepsilon=0.1$) to $2.175$ ($\varepsilon\to0$)
(Appendix~\ref{app:equilibrium}). For calibrations nearer a regime
boundary than the baseline's $8\%$ margin, the $\varepsilon$-ladder should
be re-run before any classification is reported---this is the practical
protocol the table exists to demonstrate.
\end{enumerate}

% ============================================================================
% Appendix A10 — Global Asymptotic Stability: Full Proof
% ============================================================================
\section{Global Asymptotic Stability: Full Proof of Theorem~\ref{thm:gas}}
\label{app:gas}

We assemble the complete proof from four ingredients: (1)~a forward-invariant
trapping rectangle with the \emph{corrected} ceiling $\Ctrap$;
(2)~a Bendixson--Dulac certificate excluding closed orbits;
(3)~uniqueness and local stability of the interior equilibrium; and
(4)~the boundary structure of the invariant axis $\{C=0\}$. Throughout,
parameters are positive and the interior equilibrium exists
($\Phi(O_0)>0$, Theorem~\ref{thm:uniqueness}).

\subsection{Step 1: the trapping rectangle, and why the old one leaked}

Let $\Ctrap$ be defined by Eq.~\eqref{eq:ctrap}, i.e.\ the unique capital
level at which the capital nullcline reaches the vitality ceiling,
$V_c(O(\Ctrap))=\Vmax$; at baseline
$\Ctrap=\bigl((0.1\cdot10/0.2-1)/0.15-1\bigr)/0.25=102.67$. Set
$\Omega=[0,\Vmax]\times[0,\Ctrap]$.

\emph{Forward invariance.} On the four faces of $\Omega$:
\begin{itemize}
\item $V=0$: $\dot V=R>0$ (inward; Appendix~\ref{app:invariance}).
\item $V=\Vmax$: $\dot V=-\delta O^{\gamma}\Vmax/(\Vmax+\varepsilon)<0$
(inward).
\item $C=0$: $\dot C=0$ (invariant face).
\item $C=\Ctrap$: for every $V\in[0,\Vmax]$,
\[
\dot C
= C\Bigl(\frac{\alpha V}{1+\varphi O(\Ctrap)}-\mu\Bigr)
\le C\Bigl(\frac{\alpha \Vmax}{1+\varphi O(\Ctrap)}-\mu\Bigr)
= 0,
\]
with equality only at the corner $V=\Vmax$, where the $V$-component points
strictly inward; no trajectory can exit through this face.
\end{itemize}
Moreover $\Omega$ is \emph{absorbing} for the region of interest: if
$C_0>\Ctrap$ (with $V_0\in[0,\Vmax]$), then $O>O(\Ctrap)$ implies
$\alpha V/(1+\varphi O)<\mu$, so $\dot C<0$ strictly until the trajectory
enters $\Omega$; hence every trajectory with $V_0\in[0,\Vmax]$, $C_0>0$
eventually lies in $\Omega$.

\emph{The correction.} Earlier drafts used the rectangle
$[0,\Vmax]\times[0,\Cmax]$ with the carrying capacity $\Cmax=44.99$ of
Eq.~\eqref{eq:cmax} as ceiling. That rectangle is \emph{not} forward
invariant: at $(V,C)=(\Vmax,44.99)$, $O=12.25$ and
\[
\dot C = C\Bigl(\frac{\alpha\Vmax}{1+\varphi O}-\mu\Bigr)
       = 44.99\,(0.3525-0.2) \approx +6.86 \;>\;0 ,
\]
so the flow exits through the ceiling; concretely, the baseline trajectory
from $(V_0,C_0)=(10,40)$ overshoots to $C\approx46.4$ before returning.
The carrying capacity is a flux threshold (Appendix~\ref{app:carrying});
the trapping bound is the nullcline-ceiling constant $\Ctrap$, and the two
must not be conflated. All results downstream of the trapping rectangle
are stated with $\Ctrap$; the certificate is verified numerically on
$[0,\Vmax]\times[0,\,1.2\,\Ctrap]$ by
\texttt{test\_bendixson\_dulac\_trapping\_rectangle\_is\_valid}.

\subsection{Step 2: Bendixson--Dulac certificate (no closed orbits)}

Take the Dulac function $B(V,C)=1/C$ on the open quadrant
$Q=\{V>0,\,C>0\}$, and write the field as
$\mathbf{F}=(F,G)$ with $G=C\,g(V,O)$, $g=\alpha V/(1+\varphi O)-\mu$.
Then
\[
\nabla\!\cdot\!\bigl(B\mathbf{F}\bigr)
= \frac{\partial}{\partial V}\Bigl(\frac{F}{C}\Bigr)
+ \frac{\partial}{\partial C}\bigl(g\bigr)
= \frac{1}{C}\Bigl(
   -\frac{R}{\Vmax}
   -\frac{\delta O^{\gamma}\varepsilon}{(V+\varepsilon)^{2}}
  \Bigr)
\;-\;
\frac{\alpha V\varphi\,\beta\eta C^{\eta-1}}{(1+\varphi O)^{2}} .
\]
The first bracket is strictly negative everywhere; the second term is
non-positive (strictly negative for $V>0$). Hence
$\nabla\!\cdot\!(B\mathbf{F})<0$ on all of $Q$---no sign changes, no
degenerate set. By the Bendixson--Dulac theorem
\citep{strogatz2015nonlinear}, $Q$ contains no closed orbit; the same
Green's-theorem argument excludes graphics (closed curves assembled from
orbits and equilibria) whose interior lies in $Q$. Numerically, the
supremum of the divergence over a dense grid in
$[0,\Vmax]\times(0,\,1.2\,\Ctrap]$ is $\approx-3.5\times10^{-3}$, bounded
away from zero
(\texttt{test\_bendixson\_dulac\_divergence\_strictly\_negative}).

\subsection{Step 3: equilibria inside $\Omega$}

By Theorem~\ref{thm:uniqueness} there is exactly one interior equilibrium
$(V^{*},C^{*})$; by the nullcline bound of Appendix~\ref{app:equilibrium}
it satisfies $C^{*}<\Ctrap$, so it lies in the interior of $\Omega$. By
Lemma~\ref{lem:trace} (trace negative, determinant positive) it is locally
asymptotically stable---at baseline a spiral sink
(Appendix~\ref{app:jacobian}). The only other equilibrium in $\Omega$ is
the axis point $(V_{\mathrm{axis}},0)$, $V_{\mathrm{axis}}\approx9.934$,
which is a hyperbolic saddle whose stable manifold is contained in the
invariant axis $\{C=0\}$ (Appendix~\ref{app:regularization}).

\subsection{Step 4: Poincar\'e--Bendixson assembly}

Let $(V_0,C_0)\in\Omega$ with $C_0>0$. By Step~1 the forward orbit remains
in the compact set $\Omega$, and by invariance of $\{C=0\}$ together with
uniqueness of solutions it remains in $Q$ for all finite time. Its
$\omega$-limit set $L$ is nonempty, compact, connected, and invariant. By
the Poincar\'e--Bendixson theorem, $L$ is an equilibrium, a closed orbit,
or a graphic. Step~2 excludes closed orbits and graphics in $Q$. Could $L$
meet the axis? The only invariant subsets of the axis face within $\Omega$
are the saddle and its axis-internal stable connections; if $L$ contained
the saddle, then (since $L$ is not the saddle alone unless the orbit
converges to it) $L$ would have to contain an orbit converging to the
saddle---an orbit in the saddle's stable manifold, which lies in
$\{C=0\}$ and is unreachable from $Q$. And convergence of the interior
orbit itself to the saddle is likewise impossible for the same reason.
Hence $L=\{(V^{*},C^{*})\}$: every trajectory with $C_0>0$ converges to
the interior equilibrium, which is locally stable by Step~3. This proves
global asymptotic stability on the open set $\{C>0\}\cap\Omega$, and, via
the absorbing property of Step~1, on
$\{(V,C):V\in[0,\Vmax],\,C>0\}$. \hfill$\square$

\subsection{Remarks}

\begin{enumerate}
\item \emph{The basin is open, and that is sharp.} $\{C=0\}$ is invariant:
a leader with exactly zero capital stays at zero capital, and the axis
flow converges to $(9.934,\,0)$---full vitality, no career. The theorem
therefore cannot be strengthened to the closed quadrant, and the
low-vitality equilibrium's basin is exactly $\{C>0\}$ within $\Omega$.
\item \emph{Why not a Lyapunov function.} No elementary global Lyapunov
function appears to exist: weighted quadratic forms
$a(V-V^{*})^{2}+b(C-C^{*})^{2}$ fail because the indefinite cross terms
generated by the $O^{\gamma}$ drain and the $1/(1+\varphi O)$ impact gate
cannot be dominated over all of $\Omega$ by any constant choice of
$(a,b)$. For a planar system, Dulac plus Poincar\'e--Bendixson is the
appropriate---not merely convenient---instrument; the certificate is also
mechanically checkable, which the companion tests exploit.
\item \emph{Numerical corroboration.} An $8\times8$ grid of initial
conditions spanning $[0.1,\Vmax]\times[0.5,\,2\Cmax]$ converges to
$(V^{*},C^{*})$ to within $10^{-2}$ in all 64 cases
(\texttt{test\_basin\_sweep\_all\_trajectories\_converge\_to\_zombie}),
providing an end-to-end check that no numerical artifact defeats the
analytical result.
\end{enumerate}


% ============================================================================
% BIBLIOGRAPHY
% ============================================================================
\clearpage
\bibliographystyle{apalike}
\bibliography{../references}

\end{document}
